\documentclass[11pt]{article}
\usepackage[margin=1.05in]{geometry}
\usepackage[T1]{fontenc}
\usepackage{lmodern}
\usepackage{amsmath,amssymb,amsthm,mathtools}
\usepackage{booktabs}
\usepackage{array}
\usepackage{enumitem}
\usepackage{xcolor}
\usepackage{xurl}
\usepackage[colorlinks=true,linkcolor=blue!45!black,citecolor=blue!45!black,urlcolor=blue!45!black]{hyperref}

\setlist{itemsep=2pt,topsep=4pt}

% ---------- theorem environments ----------
\newtheorem{theorem}{Theorem}[section]
\newtheorem{lemma}[theorem]{Lemma}
\newtheorem{proposition}[theorem]{Proposition}
\newtheorem{corollary}[theorem]{Corollary}
\theoremstyle{definition}
\newtheorem{definition}[theorem]{Definition}
\newtheorem{example}[theorem]{Example}
\theoremstyle{remark}
\newtheorem{remark}[theorem]{Remark}
\theoremstyle{plain}

% ---------- notation ----------
\newcommand{\UST}{\mathrm{UST}}
\newcommand{\USF}{\mathrm{USF}}
\newcommand{\TV}{\mathrm{TV}}
\newcommand{\BC}{\mathrm{BC}}
\newcommand{\E}{\mathbb{E}}
\newcommand{\R}{\mathbb{R}}
\newcommand{\eps}{\varepsilon}
\newcommand{\supp}{\operatorname{supp}}
\newcommand{\norm}[1]{\lVert #1\rVert}
\newcommand{\opnorm}[1]{\lVert #1\rVert_{\mathrm{op}}}
\newcommand{\abs}[1]{\lvert #1\rvert}
\newcommand{\cC}{\mathcal{C}}
\newcommand{\cT}{\mathcal{T}}
\newcommand{\cU}{\mathcal{U}}
\newcommand{\one}{\mathbf{1}}
\newcommand{\sqrtm}[1]{\sqrt{#1}^{\,\circ}}
\emergencystretch=3em

\title{Uniform spanning trees of expander prisms and squares\\ versus nonnegative probabilistic circuits}
\author{Chris Blaschke\\Ax0s}
\date{November 3, 2026. Preprint, not peer reviewed.}

\begin{document}
\maketitle

\begin{abstract}
Let $H$ be a connected graph with maximum degree at most $\Delta$ and edge expansion at least $h_0>0$, and let $G$ be the prism $H\square K_2$ or the square $H^2$. Let $P$ be the uniform spanning tree of $G$, viewed as a distribution on the $M$ edge indicators, and let $N$ be the number of vertices. We prove that every smooth decomposable probabilistic circuit with nonnegative weights and leaves of arity at most $2M/3$ whose distribution is within total variation $\eps<1$ of $P$ has size at least $(1-\eps)^2e^{cN}$, with an explicit $c=c(\Delta,h_0)>0$. Determinism is not assumed. By the matrix-tree theorem, $P$ has a generating circuit of size polynomial in $N$, and all its marginals are ratios of determinants. Since $M=\Theta(N)$, the optimal size of nonnegative smooth decomposable circuits for these distributions is $2^{\Theta(M)}$ at every fixed error below one.

The proof bounds, for every balanced split of the edges, the largest Bhattacharyya coefficient between $P$ and a product distribution across the split. We pass to two indicator bits per short cycle that crosses the split, and bound the square-root matrix of their joint law with an adaptive chain rule. The spatial Markov property of the uniform spanning tree gives the conditional lower bound on two local events needed for the chain rule. The information-theoretic tools are known: the chain rule is a case of an exercise in Polyanskiy and Wu, and the identification of the optimized affinity with a matrix norm is due to Lapidoth and Pfister. The circuit reduction is from our earlier paper on dimers. We believe the local cycle-bit reduction for spanning trees and the resulting constant-error bound for a natural distribution on sparse bounded-degree families are new. The constant $c$ is very small (about $2\cdot10^{-8}$ for a cubic example).
\end{abstract}

\tableofcontents

\section{Introduction}\label{sec:intro}

Probabilistic circuits (PCs) represent a distribution on $\{0,1\}^M$ by an arithmetic circuit with nonnegative weights. Smoothness and decomposability are the structural conditions under which marginals can be computed in time linear in the circuit size~\cite{ChoiVergariVdB20}. Which tractable distributions fall outside this class? We study the uniform spanning tree. By Kirchhoff's matrix-tree theorem~\cite{Kirchhoff1847} its generating polynomial is a determinant, so it has a polynomial-size generating circuit with subtraction, in the sense of Zhang, Juba and Van den Broeck~\cite{ZhangJubaVdB21}. Exact representation of uniform spanning trees of the complete graph by decomposable circuits needs exponential size~\cite{MartensMedabalimi14,ZhangHoltzenVdB20}. Those results concern exact representation or vanishing error on the complete graph, where the number of variables is quadratic in the number of vertices.

We prove lower bounds at constant total-variation error, without assuming determinism, for 2-connected bounded-degree expanders, the complete graph and dense graphs with a linear spectral gap. The two arguments behind these bounds are dual cases of one argument about matroids.

\subsection{Main result}

For a graph $G$, the edge expansion is $h(G)=\min\{\abs{\partial S}/\abs S:\ 0<\abs S\le\abs{V(G)}/2\}$, where $\partial S$ is the set of edges with exactly one endpoint in $S$. A graph is 2-connected if it has at least three vertices and stays connected after deleting any one vertex. Circuit terms are defined in Section~\ref{sec:models}. The size of a circuit is the number of nodes plus the number of wires.

\begin{theorem}[main theorem]\label{thm:main}
Fix an integer $D\ge2$ and a real $h>0$. Let $G$ be a 2-connected simple graph with $N$ vertices, $M$ edges, maximum degree at most $D$ and edge expansion $h(G)\ge h$, and let $P$ be the uniform distribution on the spanning trees of $G$, viewed as a distribution on $\{0,1\}^{E(G)}$. Let $C$ be a smooth decomposable probabilistic circuit with nonnegative weights over the $M$ edge variables, whose leaves have arity at most $2M/3$, and let $Q_C$ be its normalized distribution. If $\TV(P,Q_C)\le\eps<1$, then
\[
  \abs C\ \ge\ (1-\eps)^2\,e^{c N},\qquad
  c=c(D,h)=\frac{h_1^{5D+4}}{(D+1)^{4(D+1)}\,e^{10(D+1)}},\qquad h_1=\min(h,1).
\]
If moreover the Laplacian of $G$ has second smallest eigenvalue $\lambda_2(L_G)\ge\lambda>1$, the same holds with
\[
  c=c_\lambda(D)=\frac{\lambda\,(D-1)\,(1-1/\lambda)^D}{3D^3(D+\lambda/2)(D^2+1)} .
\]
\end{theorem}

The leaf-arity condition cannot be dropped: a single leaf of arity $M$ can store $P$. The bound $2M/3$ comes from the circuit reduction (Lemma~\ref{lem:decomp}). Girth plays no role. Theorem~\ref{thm:main} applies, with probability tending to one, to random $d$-regular graphs for every fixed $d\ge3$, to the Ramanujan graphs of Lubotzky, Phillips and Sarnak~\cite{LPS88}, and to the prisms and squares of bounded-degree expanders (Corollary~\ref{cor:families}).

\begin{corollary}[separation]\label{cor:separation}
Fix $D$, $h>0$, $\eps<1$, and a family of graphs $G$ as in Theorem~\ref{thm:main}.
\begin{enumerate}[label=(\alph*)]
\item Every nonnegative smooth decomposable PC with leaves of arity at most $2M/3$ that approximates $P$ within total variation $\eps$ has size $2^{\Omega(N)}=2^{\Omega(M)}$. Some such circuit has size $2^{O(M)}$, so $2^{\Theta(M)}$ is the optimal size.
\item The generating polynomial $\sum_xP(x)\prod_{e:x_e=1}z_e$ has an arithmetic circuit with subtraction of size polynomial in $N$, and every marginal probability $P(x_S=a)$ can be computed in polynomial time (Proposition~\ref{prop:upper}).
\end{enumerate}
\end{corollary}

For fixed $D$ the constant $c(D,h)$ is polynomial in $h$, with $\log(1/c)=O(D\log((D+1)/\min(h,1)))$; the previous version of this paper had $c\le D^{-4D/h}$, exponentially small in $D/h$. For $D=3$ and $h=1/5$ the closed form gives $c\approx5\cdot10^{-41}$ and the optimized parameters of Proposition~\ref{prop:touch-param} about $5\cdot10^{-24}$, compared with $9\cdot10^{-44}$ before (Table~\ref{tab:constants}). The spectral constant needs no small parameter at all: for 6-regular Ramanujan graphs, $\lambda_2\ge6-2\sqrt5$ gives $c\ge8.0\cdot10^{-8}$, and $c_\lambda(D)\ge2.7\cdot10^{-9}$ for all $d$-regular Ramanujan graphs with $6\le d\le64$. For Ramanujan graphs $c_\lambda(D)$ decays only polynomially in $D$, like $D^{-4}$. The degree of the polynomial in $h$ in $c(D,h)$ grows with $D$, and for the star bits behind it this is unavoidable: on 2-connected hub graphs they certify at most $O(h^{D/2+1})$ for some balanced split (Corollary~\ref{cor:hub}). A second, local argument gives much better constants on two explicit families. The prism $H\square K_2$ consists of two copies of $H$ joined by a perfect matching between corresponding vertices. The square $H^2$ has the vertices of $H$, with an edge between any two vertices at distance one or two in $H$.

\begin{corollary}[prisms and squares]\label{cor:prism-square}
Fix an integer $\Delta\ge2$ and a real $h_0>0$. Let $H$ be a connected simple graph with at least three vertices, maximum degree at most $\Delta$ and $h(H)\ge h_0$, and let $G$ be $H\square K_2$ or $H^2$. Then $G$ is 2-connected, and the conclusion of Theorem~\ref{thm:main} holds with
\[
  c=\frac{\min(1,h_0)}{24\,\Delta(\Delta+1)^5(3\Delta+4)}\quad\text{for }G=H\square K_2,\qquad
  c=\frac{8h_0}{243\,\Delta^{9}(\Delta-1)(h_0+\Delta^2)}\quad\text{for }G=H^2 .
\]
\end{corollary}

These constants come from Theorem~\ref{thm:general}, which applies to every bounded-degree expander whose vertex neighbourhoods are tied together by short cycles (Definition~\ref{def:linked}). For a cubic $H$ with $h(H)\ge1/5$, the prism gives $c\ge3\cdot10^{-7}$ (Example~\ref{ex:cubic}). The same local argument gives $e^{\Omega(L)}$ for the $L\times L$ grid (Corollary~\ref{cor:grid}), which is not an expander and is not covered by Theorem~\ref{thm:main}. All the constants are tiny, and in exact computations on small graphs the bounds we prove are far from tight (Appendix~\ref{app:checks}). The results are qualitative. For the bounded-degree expanders the exponent is linear in the number of variables; for the grid it is $\Omega(L)=\Omega(\sqrt M)$.

Expansion alone is not enough. The uniform spanning tree is the product of the uniform spanning trees of the blocks of $G$, so a balanced split along blocks has Bhattacharyya coefficient one with a product distribution. Proposition~\ref{prop:blocks} shows that such a split exists exactly when no block has more than $2M/3$ edges. Pendant edges give bounded-degree expanders of this kind (Proposition~\ref{prop:pendant}). For bounded-degree expanders this is the only obstruction: blocks inherit expansion, so if the largest block $G_1$ has $M_1>2M/3$ edges, every balanced split has $\beta^2\le e^{-\Omega(M_1-2M/3)}$, and for circuits with univariate leaves the bound $e^{\Omega(\abs{V(G_1)})}$ holds without any threshold (Theorem~\ref{thm:blocks}). 2-connectivity is the case $G_1=G$.

We also record exact-support consequences (Corollary~\ref{cor:support}). For every graph as in Theorem~\ref{thm:main}, every nonnegative smooth decomposable PC with the same leaf-arity condition, and every monotone arithmetic circuit, whose support is exactly the set of spanning trees of $G$, with arbitrary positive values on that set, has size $e^{\Omega(N)}$.

\subsection{The complete graph}

Both arguments behind Theorem~\ref{thm:main} need bounded degree. Let $\mu_n$ be the uniform spanning tree of $K_n$, a distribution on $M=\binom n2$ variables. Section~\ref{sec:complete} proves that for $n\ge3$ every balanced split $A|B$ of $E(K_n)$ satisfies
\[
  -\ln\beta(\mu_n;A|B)^2\ \ge\ \frac n{146}-\frac12\ln n-3
\]
(Theorem~\ref{thm:complete}), with an explicit sum $\Phi_n$ behind it whose exponential rate is $1/144.86$ per vertex. Consequently every smooth decomposable PC with nonnegative weights and leaves of arity at most $2M/3$ with $1-\TV(\mu_n,Q_C)\ge\delta>0$ has size at least $\delta^2e^{n/146}/(e^3\sqrt n)$ (Corollary~\ref{cor:complete}). The arity condition cannot simply be dropped, since a single leaf of arity $M$ represents $\mu_n$. A subset recursion gives an exact deterministic circuit of size $O(n^43^n)$ (Proposition~\ref{prop:complete-upper}), so at every fixed error the optimal size is $2^{\Theta(n)}=2^{\Theta(\sqrt M)}$ (Corollary~\ref{cor:complete-order}). Martens and Medabalimi~\cite{MartensMedabalimi14} proved a lower bound of the same order for exact representation and for approximation arbitrarily well; ours holds at every fixed total-variation error below one. The proved rate is still far from the truth. Exact computation for $6\le n\le9$ gives $-\ln\beta^2$ between $0.12n$ and $0.30n$ on the splits we tried, and a Monte Carlo estimate, which is an upper bound on the true rate and is not certified, gives about $0.17n$ for large $n$ (Remark~\ref{rem:kn-ceiling}).

The same method gives $2^{\Omega(n)}$ for every graph whose Laplacian has $\lambda_2\ge\kappa n$, with a rate of order $\kappa^4e^{-1.14/\kappa}$ (Theorem~\ref{thm:dense} and Corollary~\ref{cor:dense}). For $G(n,\frac12)$ this applies conditionally on the standard concentration of $\lambda_2$, which we do not re-derive. A minimum degree of $n/2$ is not enough for the method: two cliques joined by a perfect matching have a balanced split with $\beta^2>e^{-2}$ (Proposition~\ref{prop:two-cliques}).

\subsection{Matroids}

A spanning tree contains no cycle and meets every cut. The local argument uses the first fact and the vertex-touch argument the second. Section~\ref{sec:matroid} shows that the two arguments are the circuit and cocircuit cases of one argument for the uniform distribution on the bases of a matroid. Each element of a circuit is certified by a cut that meets the circuit in exactly two elements, and each element of a cocircuit by a cycle; an exchange argument turns such a frame into a conditional lower bound (Lemma~\ref{lem:mat-frame}). This gives a bound for every matroid in terms of compatible framed circuits and cocircuits (Theorems~\ref{thm:mat-local} and~\ref{thm:mat-pc}), and the bound is invariant under matroid duality (Proposition~\ref{prop:mat-dual}). For balanced matroids, which include the regular ones, a global version replaces frames by leverages (Theorem~\ref{thm:mat-global}). For a totally unimodular representation the leverages are effective resistances (Lemma~\ref{lem:mat-resist}), and applied to vertex stars this recovers Theorem~\ref{thm:touch-cut}. As an example outside graphs and their duals we obtain $2^{\Omega(n)}$ bounds for binary non-regular matroids obtained from linked expanders by attaching a Fano plane to every edge (Proposition~\ref{prop:fano}). The key lemma of the global version fails for the non-balanced binary matroid $S_8$ (Section~\ref{sec:limits}); whether the global version holds for all matroids is open.

\subsection{Subtraction}

The uniform spanning tree has a small circuit with subtraction. Section~\ref{sec:signed} asks how much of the lower bound survives when subtraction is allowed but the circuit keeps a fixed decomposition structure. For structured-decomposable circuits with real weights, a rank argument adapted from de Colnet and Mengel~\cite{DeColnetMengel21} gives exponential size for exact representation on linked expanders (Theorem~\ref{thm:signed-rank} and Corollary~\ref{cor:signed-exact}). One inequality for disjoint permutation blocks of $\UST$ (Theorem~\ref{thm:block-approx}) extends this to approximation. If the output $Q$ sums to $1$, every fixed error $\frac12\norm{\UST-Q}_1\le\eps<1$ needs size $e^{\Omega(N)}$, although $Q$ may be negative (Theorem~\ref{thm:signed-o2} and Corollary~\ref{cor:signed-approx}); the key step is that a matrix of rank $r$ and bounded $\ell_1$ mass covers at most a $\sqrt{r/n}$ fraction of an $n\times n$ permutation, a consequence of a rank bound of Ky Fan and Hoffman~\cite{FanHoffman53}. Without normalization the threshold is $\eps<\frac12$, which is sharp because $Q=0$ has error $\frac12$ (Theorem~\ref{thm:signed-approx}). The same bounds hold on every 2-connected bounded-degree expander, at any girth and without smoothness (Corollary~\ref{cor:signed-expander}), with blocks built from pendant vertices instead of short cycles, and they cover sparse and low-degree approximations (Corollary~\ref{cor:signed-sparse}). Normalized squares and positive semidefinite forms of structured-decomposable amplitudes obey the same lower bound as nonnegative circuits (Theorem~\ref{thm:signed-square}): for structured circuits the barrier is decomposability, not sign. For unstructured decomposable circuits with subtraction we have lower bounds only for restricted classes. On bounded-degree expanders, support-exact representation needs size $2^{N^{\Omega(1/d)}}$ at product-depth $d$ and $N^{\Omega(\log N)}$ for formulas, and sums of $r$ structured circuits, each on its own vtree, need total size $2^{\Omega(3^{-r}\sqrt N)}$ (Corollary~\ref{cor:signed-unstructured}). These proofs transfer the restriction arguments of Raz~\cite{Raz09} and of Raz and Yehudayoff~\cite{RazYehudayoff09} through a large complete minor~\cite{KrivelevichNenadov21,Krivelevich19}, after checking that their rank bounds do not depend on the constants of the restriction. For sums of $r$ structured circuits, a packing of far-apart pendant gadgets gives total size $2^{\Omega_r(N)}$ directly in the graph, both for support-exact representation and at error $\eps<\frac12$ (Theorem~\ref{thm:signed-sums-linear}); its constants are astronomically small, so the minor bound remains better at any realistic size. General circuits stay open. For exact representation the general question is equivalent to the syntactically multilinear complexity of the spanning-tree polynomial (Proposition~\ref{prop:signed-multilinear}), and a small approximating circuit would have to cancel all but an exponentially small fraction of its absolute mass (Theorem~\ref{thm:signed-cancel}), largely between terms on different splits (Proposition~\ref{prop:signed-cross}). We explain why the block method does not yet reach unstructured circuits (Section~\ref{sec:signed-barriers}).

\subsection{What is known and what is new}\label{sec:intro-priority}

The ingredients of the proof are known; we credit them where they are used.
\begin{itemize}
\item For a joint distribution $\mu$ of $(X,Y)$, the largest Bhattacharyya coefficient between $\mu$ and a product distribution equals the top singular value of the entrywise square root of $\mu$. This is Lemma~6 of Lapidoth and Pfister~\cite{LapidothPfister19}, stated for their order-$\frac12$ dependence measure $J_{1/2}$.
\item The adaptive chain rule we use (Lemma~\ref{lem:chain}) is a case of Exercise~I.42(b) in Polyanskiy and Wu~\cite{PolyanskiyWu25}, combined with the previous item. For independent pairs it is the additivity of $J_{1/2}$~\cite[Lemma~12]{LapidothPfister19}. Hellinger distance and its product structure are standard tools for communication lower bounds~\cite{BJKS04,Jayram09}.
\item The reduction from circuits to products across balanced edge splits is the balanced decomposition and the first step of the overlap lemma from our earlier paper on dimers~\cite[Lemmas~7.1 and~7.2]{Blaschke26Dimers}. The decomposition itself is the standard balancing argument for syntactically multilinear circuits~\cite{RazYehudayoff08}, also stated for sum-product networks by Martens and Medabalimi~\cite[Thm.~38]{MartensMedabalimi14}.
\item Kirchhoff's formula for edge marginals, Rayleigh monotonicity, Thomson's principle and the triangle inequality for effective resistance are classical~\cite{Kirchhoff1847,LyonsPeres16,KleinRandic93}, as is Mohar's Cheeger-type inequality~\cite{Mohar89}, whose proof we repeat for the lower bound we use (Lemma~\ref{lem:cheeger}). The count of connected sets (Lemma~\ref{lem:count}) is a form of a lemma of Fomin and Villanger~\cite{FominVillanger12}, and the random deletion (Lemma~\ref{lem:delete}) is the alteration method with a Peierls-type union bound. Lemma~\ref{lem:forest-mgf} and Section~\ref{sec:edge-pairs} use the transfer-current theorem of Burton and Pemantle~\cite{BurtonPemantle93}. The proof of Theorem~\ref{thm:complete} uses de Caen's bound on sums of squared degrees~\cite{deCaen98}, the exact law of the isolated vertices of a uniform labelled forest (a Cayley-type count), and Robbins' form of Stirling's formula~\cite{Robbins55}.
\item For matroids we use standard facts from Oxley~\cite{Oxley11}, the negative correlation of regular matroids due to Feder and Mihail~\cite{FederMihail92}, and the characterization of balanced binary matroids by Choe and Wagner~\cite{ChoeWagner03}.
\end{itemize}
Constant-error lower bounds against nonnegative decomposable circuits are not new in general. Leland and Choi~\cite{LelandChoi25} prove one for \emph{deterministic} circuits. A generic $\exp(\Omega(n))$ separation from a random linear code, whose inference is easy by Gaussian elimination, follows from the rank bounds of de Colnet and Mengel~\cite{DeColnetMengel20}, as explained in~\cite[Sec.~10.4]{Blaschke26Dimers}. The dimer paper~\cite{Blaschke26Dimers} proves an $\exp(\Theta(\sqrt n))$ bound for the uniform domino tiling of the square grid. That paper also lists the robust spanning-tree question for expanders as open~\cite[Sec.~11, problem~5]{Blaschke26Dimers}.

We believe the following are new.
\begin{enumerate}[label=(\roman*)]
\item The constant-error $2^{\Omega(N)}$ lower bound against nonnegative smooth decomposable PCs, without determinism, for uniform spanning trees of every 2-connected bounded-degree expander, at any girth. Apart from an earlier version of this paper, which covered prisms and squares of expanders, we know of no constant-error exponential lower bound against nondeterministic decomposable PCs for uniform spanning trees of any family, and none for a distribution of this kind whose inference reduces to determinants.
\item The vertex-touch reduction: two bits per vertex that has edges on both sides of the split, conditioning only on the stars of a sparse set of such vertices, and a conditional lower bound from effective resistances, in a telescoped form controlled by $\lambda_2$ of the graph left after the sparse set is deleted, with a random deletion that keeps this graph an expander (Lemma~\ref{lem:touch}, Theorem~\ref{thm:touch-cut}, Lemmas~\ref{lem:touch-flow} and~\ref{lem:delete}).
\item The local cycle-bit reduction: for spanning trees, two indicator bits per crossing cycle, together with the spatial Markov property at the edges incident to the cycle, give the conditional non-degeneracy that the chain rule needs (Lemma~\ref{lem:local} and Theorem~\ref{thm:cut}). It gives the better constants of Corollary~\ref{cor:prism-square} and the grid bound.
\item The constant-error $2^{\Omega(n)}$ bound for the uniform spanning tree of $K_n$ with an explicit rate (Theorem~\ref{thm:complete} and Corollary~\ref{cor:complete}), its extension to dense graphs (Corollary~\ref{cor:dense}), and the averaged conditioning on a rectangle (Lemma~\ref{lem:avg-cond}) that lets the proof use a conditioning set that is not independent. We did not find Lemma~\ref{lem:avg-cond} stated in this form, but it is elementary.
\item The matroid form of both arguments: framed circuits and cocircuits, the local theorem for every matroid, the global theorem for balanced matroids, and the Fano-decorated family (Section~\ref{sec:matroid}). We did not find the exchange bound of Lemma~\ref{lem:mat-exchange} stated for matroids, but it is also elementary.
\item The approximate lower bounds for structured-decomposable circuits with real weights: the permutation-block inequality (Theorem~\ref{thm:block-approx}), which reaches every error below $1$ for normalized outputs, and its pendant-vertex form on every 2-connected bounded-degree expander, without smoothness (Corollary~\ref{cor:signed-expander}). Lemma~\ref{lem:kfh} is classical~\cite{FanHoffman53}.
\item The support-exact lower bounds for restricted unstructured signed circuits (Corollary~\ref{cor:signed-unstructured}), and the bounds for sums of structured circuits through packings of pendant gadgets (Theorem~\ref{thm:signed-sums-linear}). We did not find them in the literature, but the transfer is short and may be folklore; the rank arguments are those of~\cite{Raz09,RazYehudayoff09}.
\item The limits of the vertex-touch method: on 2-connected hub graphs with expansion $h$ and maximum degree $D$, star bits certify at most $O(h^{D/2+1})$ for one balanced split, while bits for pairs of edges at a vertex certify $h^3/(10935D^5)$ for that split (Corollary~\ref{cor:hub} and Proposition~\ref{prop:hub-pairs}). No novelty is claimed for the ingredients.
\end{enumerate}
Our literature search was targeted, not exhaustive. The closest results we found are discussed in Section~\ref{sec:related}. In particular, Chattopadhyay, Datta and Mukhopadhyay~\cite{CDM20Negations} and Chattopadhyay, Datta, Ghosal and Mukhopadhyay~\cite{CDGM22} prove exact $2^{\Omega(n)}$ monotone lower bounds for directed spanning-tree polynomials of expanders of large constant degree, and M\'esz\'aros and Peled~\cite{MeszarosPeled26} prove that the uniform spanning tree of a graph with large minimum degree is close to a product-like process in expected edit distance. Neither statement concerns total variation for circuits; Section~\ref{sec:related} explains why neither contradicts nor implies ours.

\subsection{Proof outline}\label{sec:outline}

Fix a split $E=A\sqcup B$ of the edges. For a product distribution $q=q^A\otimes q^B$ write $\BC(P,q)=\sum_x\sqrt{P(x)q(x)}$, and let $\beta(P;A|B)$ be the largest value over all products. The circuit reduction (Section~\ref{sec:circuits}) gives
\[
  \abs C\ \ge\ (1-\eps)^2\Big/\max_{A|B\ \text{balanced}}\beta(P;A|B)^2 .
\]
So it suffices to show $\beta\le e^{-cN/2}$ for every split in which both sides have at least $M/3$ edges. The proof of Theorem~\ref{thm:main} has four steps.
\begin{enumerate}
\item \emph{Vertex-touch bits.} Call a vertex $w$ mixed if it has edges in both $A$ and $B$. Put $a_w=\one[T\text{ uses no }A\text{-edge at }w]$ and $b_w=\one[T\text{ uses no }B\text{-edge at }w]$. The bits $a$ depend only on $T\cap A$ and the bits $b$ only on $T\cap B$. A spanning tree has an edge at every vertex, so $a_w=b_w=1$ never happens. Data processing gives $\beta(P;A|B)\le\opnorm{\sqrt\mu}$, where $\mu$ is the joint law of the bits at a set $W$ of mixed vertices (Lemma~\ref{lem:dpi}).
\item \emph{Chain rule.} If, given every possible value of the earlier pairs, the pair $(a_w,b_w)$ has probability at least $p_A$ of being $(1,0)$ and at least $p_B$ of being $(0,1)$, then the pair contributes a factor $((1+\sqrt{1-4p_Ap_B})/2)^{1/2}\le e^{-p_Ap_B/2}$ to $\opnorm{\sqrt\mu}$ (Lemmas~\ref{lem:chain} and~\ref{lem:two-by-two}). No independence between the pairs is needed.
\item \emph{Conditioning on stars.} Take $W$ independent and condition on the tree on the stars of the other vertices of $W$. The rest of the tree is a uniform spanning tree of a graph obtained from $G-(W\setminus\{w\})$ by identifying vertices. Deleting the $A$-edges at $w$ one at a time, Kirchhoff's formula and Rayleigh monotonicity give $\Pr[a_w=1\,|\,\text{past}]\ge\prod_u(1+R_u)^{-1}$, with effective resistances $R_u$ in a graph that contains $G-W$ (Lemma~\ref{lem:touch}). Spreading the current from $u$ over all remaining neighbours of $w$, the product telescopes to $\Pr[a_w=1\,|\,\text{past}]\ge(b/d)(1+1/\lambda_2(L_{G-W}))^{-a}$, where $a$ and $b$ are the numbers of $A$- and $B$-edges at $w$ and $d=a+b$ (Lemma~\ref{lem:touch-flow}).
\item \emph{Expansion.} Every balanced split of an expander has linearly many mixed vertices (Lemma~\ref{lem:mixed-vertices}). If $G$ is 2-connected, a random sparse set $W$ of them, with the vertices of a few dangerous configurations removed, leaves $G-W$ an expander (Lemma~\ref{lem:delete}); the dangerous configurations are counted by a Peierls argument over connected sets (Lemma~\ref{lem:count}). Mohar's inequality turns the expansion of $G-W$ into a bound on $\lambda_2$ (Lemma~\ref{lem:cheeger}). If $\lambda_2(L_G)>1$, a set $W$ with pairwise distances at least $3$ suffices, since deleting it lowers $\lambda_2$ by at most one (Lemma~\ref{lem:sep3}).
\end{enumerate}
The local argument for Corollary~\ref{cor:prism-square} replaces step~1 by two bits per short cycle that crosses the split ($u=\one[\text{all }A\text{-edges of the cycle are in }T]$ and $v$ likewise for $B$; a tree contains no cycle), step~3 by conditioning on everything outside the edges incident to the cycle and counting completions (Lemma~\ref{lem:local}), and step~4 by a packing of short cycles (Lemma~\ref{lem:mixed}). Cycles and cuts give the two arguments: a tree contains no cycle and meets every cut. Section~\ref{sec:matroid} turns this remark into a theorem about matroids. Only steps~1 and~3 use spanning trees.

For $K_n$ (Section~\ref{sec:complete}) the set $W$ cannot be independent. We condition on the tree edges inside $W$, keep the set $I$ of vertices of $W$ that these edges do not touch, and average the resulting bounds over the conditioning (Lemma~\ref{lem:avg-cond}). Given the conditioning, a star-deletion bound that is sharp on complete graphs gives the kernel of each vertex of $I$ (Lemma~\ref{lem:kn-star}). The law of $I$ for a uniform $r$-set $W$ is explicit (Lemma~\ref{lem:kn-iso}), de Caen's inequality shows that an average vertex has a constant fraction of its edges on each side, and choosing $r\approx0.38n$ gives the rate. For dense graphs, the transfer-current theorem and Hadamard's inequality replace the exact law of $I$ (Lemma~\ref{lem:forest-mgf}). The determinant enters only on the upper-bound side.

\paragraph{Organization.} Sections~\ref{sec:prelim} and~\ref{sec:tools} give notation, circuit models and the known tools. Section~\ref{sec:vertex} proves the bound for 2-connected expanders, and Section~\ref{sec:complete} adapts it to the complete graph. Sections~\ref{sec:trees} and~\ref{sec:geometry} give the local cycle-bit argument. Section~\ref{sec:main} proves the main results, with constants, the families covered, the upper bound and exact-support consequences. Section~\ref{sec:matroid} puts both arguments in the setting of matroids. Section~\ref{sec:signed} treats circuits with subtraction. Section~\ref{sec:limits} discusses limitations and variants, including the block structure and the matroid $S_8$; Sections~\ref{sec:related} and~\ref{sec:open} cover related work and open questions. Appendix~\ref{app:checks} describes the machine checks.

\section{Preliminaries}\label{sec:prelim}

\paragraph{Graphs.} All graphs are finite. $G=(V,E)$ is connected and simple unless we say otherwise, with $N:=\abs V\ge2$, $M:=\abs E$ and maximum degree $D$. For $S\subseteq V$, $\partial S$ is the set of edges with exactly one endpoint in $S$. The edge expansion is
\[
  h(G):=\min\Bigl\{\frac{\abs{\partial S}}{\abs S}:\ S\subseteq V,\ 0<\abs S\le N/2\Bigr\}.
\]
Taking $\abs S=1$ shows $h(G)\le D$. A cycle $C$ is a simple cycle with vertex set $V(C)$ and edge set $E(C)$, and $\abs C:=\abs{E(C)}=\abs{V(C)}\ge3$. Two cycles are vertex-disjoint if their vertex sets are disjoint. For a cycle $C$, let
\[
  I(C):=\{e\in E:\ e\text{ has an endpoint in }V(C)\}
\]
be the set of edges incident to $C$. It contains $E(C)$ and every chord of $C$.

\paragraph{Spanning trees.} Let $\cT(G)$ be the set of spanning trees and $\tau(G):=\abs{\cT(G)}$. We identify an edge set with its indicator vector in $\{0,1\}^E$. The uniform spanning tree $\UST(G)$ is the uniform distribution on $\cT(G)$, viewed as a distribution on $\{0,1\}^E$. We write $T$ for a random tree with this law.

\paragraph{Splits and products.} An \emph{edge split} is a partition $E=A\sqcup B$. It is \emph{balanced} if $\abs A\ge M/3$ and $\abs B\ge M/3$. We write $x=(x_A,x_B)$ for $x\in\{0,1\}^E$. A \emph{product distribution across $A|B$} is $q=q^A\otimes q^B$, where $q^A$ and $q^B$ are probability distributions on $\{0,1\}^A$ and $\{0,1\}^B$. The product may put mass on non-trees.

\paragraph{Distances.} For distributions $P,Q$ on a finite set, $\TV(P,Q)=\frac12\sum_x\abs{P(x)-Q(x)}$, so $1-\TV(P,Q)=\sum_x\min(P(x),Q(x))$. The Bhattacharyya coefficient is $\BC(P,Q)=\sum_x\sqrt{P(x)Q(x)}$. In the normalization of Polyanskiy and Wu, $H^2(P,Q)=2-2\BC(P,Q)$, and the R\'enyi divergence of order $\frac12$ is $D_{1/2}(P\|Q)=-2\log\BC(P,Q)$. The central quantity of this paper is
\[
  \beta(P;A|B):=\max\bigl\{\BC(P,q^A\otimes q^B)\bigr\},
\]
the maximum over all product distributions across $A|B$.

\paragraph{Matrices.} For a nonnegative function $\mu$ on a product $\mathcal X\times\mathcal Y$ of finite sets, $\sqrt\mu$ denotes the $\mathcal X\times\mathcal Y$ matrix with entries $\sqrt{\mu(x,y)}$ (entrywise square root). $\opnorm W$ is the largest singular value of a matrix $W$.

\subsection{Circuit models}\label{sec:models}

We use the definitions of~\cite[Sec.~2.1]{Blaschke26Dimers}, which we repeat.

\paragraph{Probabilistic circuits.} A \emph{probabilistic circuit} (PC) over $x\in\{0,1\}^M$ is a directed acyclic graph with one root and three kinds of node.
\begin{itemize}
\item \emph{Leaves} are nonnegative functions of at most $a$ of the variables. The \emph{arity} is $a$. Indicator leaves $[x_e=0]$ and $[x_e=1]$ are the usual choice.
\item \emph{Sum nodes} compute a weighted sum of their children with nonnegative weights.
\item \emph{Product nodes} multiply their children.
\end{itemize}
The \emph{scope} $X_v$ of a node $v$ is the set of variables appearing in leaves below $v$, and $v$ computes a function $f_v$ of $x_{X_v}$. The circuit computes $f_C:=f_{\mathrm{root}}\ge0$ and represents $Q_C:=f_C/\sum_xf_C(x)$. It is \emph{decomposable} if the children of each product node have pairwise disjoint scopes, and \emph{smooth} if the children of each sum node have equal scopes. It is \emph{deterministic} if the children of each sum node have pairwise disjoint supports. We never assume determinism.

\paragraph{Size convention.} The size $\abs C$ is the number of nodes plus the number of wires. This is the convention under which the decomposition lemma (Lemma~\ref{lem:decomp}) gives at most $\abs C$ terms. For circuits of bounded fan-in, it differs from the node count by a constant factor. Remark~\ref{rem:nodes} explains how to obtain the same bound with nodes only.

\paragraph{Arithmetic circuits.} A \emph{monotone arithmetic circuit} has binary $+$ and $\times$ gates, inputs $z_e$ ($e\in E$) and nonnegative real constants. It computes a formal polynomial, in which $z_e^2\ne z_e$. Its size is its number of nodes. The \emph{generating polynomial} of a distribution $Q$ on $\{0,1\}^E$ is $\sum_xQ(x)z^x$, where $z^x:=\prod_{e:x_e=1}z_e$. A \emph{probabilistic generating circuit}~\cite{ZhangJubaVdB21} is an arithmetic circuit, with subtraction allowed, that computes it.

\section{Known tools}\label{sec:tools}

This section collects the tools the proof uses. None of them is new. We give short proofs where this keeps the paper self-contained, and say where each statement comes from.

\subsection{Affinity with products is a matrix norm}

Let $\mu$ be a joint distribution of $(X,Y)$ on a finite product $\mathcal X\times\mathcal Y$. Write $\beta(\mu;X|Y)$ for the largest value of $\BC(\mu,a\otimes b)$ over probability distributions $a$ on $\mathcal X$ and $b$ on $\mathcal Y$. For a distribution $P$ on $\{0,1\}^E$ and a split $A|B$ we view $P$ as a joint law of $(x_A,x_B)$, so that $\beta(P;A|B)$ is the quantity defined in Section~\ref{sec:prelim}.

\begin{lemma}[Lapidoth and Pfister~{\cite[Lemma~6]{LapidothPfister19}}]\label{lem:LP}
$\beta(\mu;X|Y)=\opnorm{\sqrt\mu}$, where $\sqrt\mu$ is the $\mathcal X\times\mathcal Y$ matrix with entries $\sqrt{\mu(x,y)}$.
\end{lemma}

\begin{proof}
For product $a\otimes b$, put $f=\sqrt a$ and $g=\sqrt b$. These are unit vectors, and $\BC(\mu,a\otimes b)=f^{\top}\sqrt\mu\,g\le\opnorm{\sqrt\mu}$. Conversely, $\opnorm{\sqrt\mu}=\max f^{\top}\sqrt\mu\,g$ over unit vectors, and since $\sqrt\mu$ has nonnegative entries the maximum is attained at nonnegative $f,g$. Then $a=f^2$ and $b=g^2$ are probability distributions with $\BC(\mu,a\otimes b)=\opnorm{\sqrt\mu}$.
\end{proof}

In the notation of~\cite{LapidothPfister19}, $J_{1/2}(X;Y)=-2\log\opnorm{\sqrt\mu}$ is the R\'enyi divergence of order $\frac12$ between $\mu$ and the closest product distribution. The rows and columns of $\sqrt\mu$ are indexed by the whole values of $X$ and $Y$. No independence between coordinates is assumed anywhere.

\subsection{Data processing}

\begin{lemma}[data processing]\label{lem:dpi}
Let $\phi:\mathcal X\to\mathcal X'$ and $\psi:\mathcal Y\to\mathcal Y'$ be maps, and let $\mu'$ be the law of $(\phi(X),\psi(Y))$. Then $\beta(\mu;X|Y)\le\beta(\mu';\phi(X)|\psi(Y))=\opnorm{\sqrt{\mu'}}$.
\end{lemma}

\begin{proof}
Let $a\otimes b$ be a product distribution, and let $a'$ and $b'$ be the laws of $\phi$ under $a$ and of $\psi$ under $b$. The fibre of $(x',y')$ is $\phi^{-1}(x')\times\psi^{-1}(y')$, and $\sum a(x)b(y)$ over this fibre is $a'(x')b'(y')$. By the Cauchy--Schwarz inequality on each fibre,
\[
  \BC(\mu,a\otimes b)=\sum_{x',y'}\ \sum_{\substack{x\in\phi^{-1}(x')\\ y\in\psi^{-1}(y')}}\sqrt{\mu(x,y)}\sqrt{a(x)b(y)}
  \le\sum_{x',y'}\sqrt{\mu'(x',y')\,a'(x')b'(y')}=\BC(\mu',a'\otimes b').
\]
The right side is at most $\beta(\mu')$, and Lemma~\ref{lem:LP} identifies $\beta(\mu')$ with the norm.
\end{proof}

This is the data-processing inequality for $J_{1/2}$, applied on both sides~\cite[Lemma~4]{LapidothPfister19}. It is the step that lets us discard almost all of the tree. Coarsening can only increase the affinity with products, so it is the right direction for an upper bound on $\beta$. The two maps must act separately on the two sides; this is what keeps the image of a product distribution a product.

\subsection{An adaptive chain rule}

Let $(X_k,Y_k)_{k=1}^r$ be random variables with values in finite sets $\mathcal X_k\times\mathcal Y_k$ and joint law $\mu$. Put $X=(X_1,\dots,X_r)$ and $Y=(Y_1,\dots,Y_r)$, and view $\sqrt\mu$ as a matrix with rows indexed by values of $X$ and columns by values of $Y$. For each $k$ and each value $\pi$ of the joint past $(X_{<k},Y_{<k})$ that has positive probability, let $K_k(\cdot\,|\,\pi)$ be the conditional law of $(X_k,Y_k)$, viewed as an $\mathcal X_k\times\mathcal Y_k$ matrix, and put
\[
  \gamma_k:=\max_\pi\ \opnorm{\sqrt{K_k(\cdot\,|\,\pi)}} .
\]

\begin{lemma}[adaptive chain rule]\label{lem:chain}
$\opnorm{\sqrt\mu}\le\gamma_1\gamma_2\cdots\gamma_r$.
\end{lemma}

\begin{proof}
Induction on $r$. For $r=1$ the past is empty and $\gamma_1=\opnorm{\sqrt\mu}$. Let $r\ge2$, write $x=(x',x_r)$ and $y=(y',y_r)$, and let $\mu'$ be the law of the first $r-1$ pairs. Then $\mu(x,y)=\mu'(x',y')K_r(x_r,y_r\,|\,x',y')$ when $\mu'(x',y')>0$, and $\mu(x,y)=0$ otherwise. Let $f,g$ be unit vectors. For each $(x',y')$ with $\mu'(x',y')>0$, apply the definition of $\gamma_r$ to the vectors $f(x',\cdot)$ and $g(y',\cdot)$:
\begin{align*}
  \sum_{x,y}\sqrt{\mu(x,y)}f(x)g(y)
  &=\sum_{x',y'}\sqrt{\mu'(x',y')}\sum_{x_r,y_r}\sqrt{K_r(x_r,y_r\,|\,x',y')}\,f(x',x_r)\,g(y',y_r)\\
  &\le\gamma_r\sum_{x',y'}\sqrt{\mu'(x',y')}\,F(x')\,G(y'),
\end{align*}
where $F(x')=\norm{f(x',\cdot)}_2$ and $G(y')=\norm{g(y',\cdot)}_2$. Since $\norm F_2=\norm f_2=1$ and $\norm G_2=1$, the right side is at most $\gamma_r\opnorm{\sqrt{\mu'}}$. The conditional laws of $\mu'$ are those of $\mu$ for $k<r$, so induction gives $\opnorm{\sqrt{\mu'}}\le\gamma_1\cdots\gamma_{r-1}$.
\end{proof}

The key point is that $F$ depends only on $x'$ and $G$ only on $y'$, even though the kernel $K_r$ depends on both. By Lemma~\ref{lem:LP}, Lemma~\ref{lem:chain} says $J_{1/2}(X;Y)\ge\sum_k\min_\pi J_{1/2}(X_k;Y_k\,|\,\pi)$, where the minimum is over joint pasts. This is a special case of a known chain rule. Polyanskiy and Wu~\cite[Exercise~I.42(b)]{PolyanskiyWu25} ask the reader to show that
\[
  \BC(P_{Z^n},Q_{Z^n})\le\prod_{i=1}^n\sup_{h}\BC\bigl(P_{Z_i|Z^{i-1}=h},\,Q_{Z_i|Z^{i-1}=h}\bigr)
\]
for any two joint laws (they write $1-\frac12H^2$ for $\BC$). Take $Z_i=(X_i,Y_i)$, $P=\mu$, and $Q=a\otimes b$ a product across $X|Y$. Each conditional $Q_{Z_i|Z^{i-1}=h}$ is then a product across $X_i|Y_i$, so each factor is at most $\gamma_i$ by Lemma~\ref{lem:LP}. Maximizing over $Q$ gives Lemma~\ref{lem:chain}. For independent pairs the lemma is the additivity of $J_{1/2}$~\cite[Lemma~12]{LapidothPfister19}, or the multiplicativity of operator norms under tensor products.

\begin{remark}[what ``adaptive'' means]\label{rem:adaptive}
In Lemma~\ref{lem:chain} the kernels may depend on the entire joint past of both coordinates. The coordinates themselves, and their order, are fixed in advance. In our application the cycles and their order are chosen before the tree is sampled. Choosing the next statistic as a function of both parties' earlier bits would not be covered: such a choice is not a function of one side alone, so Lemma~\ref{lem:dpi} would not apply.
\end{remark}

The conditional hypothesis cannot be replaced by a hypothesis on marginals. Let $(X_1,Y_1)$ be $(0,1)$ or $(1,0)$ with probability $\frac12$ each, and copy it to every later pair. Each pair alone has norm $1/\sqrt2$, but $\opnorm{\sqrt\mu}=1/\sqrt2$ for every $r$. Lemma~\ref{lem:chain} is consistent with this, since after the first pair the kernels are point masses and $\gamma_k=1$.

\begin{lemma}[one forbidden cell]\label{lem:two-by-two}
Let $K$ be a distribution on $\{0,1\}^2$ with $K(1,1)=0$, $K(0,1)\ge p$ and $K(1,0)\ge p$. Then
\[
  \opnorm{\sqrt K}^2=\frac{1+\sqrt{1-4K(0,1)K(1,0)}}2\le\frac{1+\sqrt{1-4p^2}}2\le1-p^2 .
\]
\end{lemma}

\begin{proof}
Write $a=K(0,0)$, $b=K(0,1)$, $c=K(1,0)$. The singular values $\sigma_1\ge\sigma_2$ of $\sqrt K=\left(\begin{smallmatrix}\sqrt a&\sqrt b\\ \sqrt c&0\end{smallmatrix}\right)$ satisfy $\sigma_1^2+\sigma_2^2=a+b+c=1$ and $\sigma_1\sigma_2=\abs{\det\sqrt K}=\sqrt{bc}$. So $\sigma_1^2$ is the larger root of $s^2-s+bc=0$, which is the stated formula, and it decreases in $bc\ge p^2$. The last step is $\sqrt{1-x}\le1-x/2$.
\end{proof}

The middle bound is attained by independent copies of the kernel with $K(0,1)=K(1,0)=p$, so for kernels with one forbidden cell, Lemmas~\ref{lem:chain} and~\ref{lem:two-by-two} lose nothing.

\subsection{From circuits to balanced products}\label{sec:circuits}

The circuit reduction is taken from our dimer paper~\cite[Sec.~7.1--7.2]{Blaschke26Dimers}. Only its first, distribution-free part is needed here. The matching-profile part of that paper is specific to perfect matchings and plays no role.

\begin{lemma}[overlap with sums of products; first step of {\cite[Lemma~7.1]{Blaschke26Dimers}}]\label{lem:overlap}
Let $P$ be any distribution on $\{0,1\}^E$, and let
\[
  Q=\sum_{u=1}^kg_u(x_{A_u})\,h_u(x_{B_u}),
\]
where $g_u,h_u\ge0$, each $A_u\sqcup B_u=E$ is a balanced split (the splits may differ between terms), and $\sum_xQ(x)=1$. Then
\[
  1-\TV(P,Q)\le\sqrt k\ \max_{A|B\ \mathrm{balanced}}\beta(P;A|B).
\]
\end{lemma}

\begin{proof}
Drop zero terms and write $g_uh_u=\lambda_uq_u$, where $q_u$ is a product distribution across $A_u|B_u$, $\lambda_u\ge0$ and $\sum_u\lambda_u=1$. Pointwise, $\min(P,Q)\le\sqrt{PQ}$ and $\sqrt{\sum_u\lambda_uq_u}\le\sum_u\sqrt{\lambda_uq_u}$. Hence
\[
  1-\TV(P,Q)=\sum_x\min(P(x),Q(x))\le\sum_u\sqrt{\lambda_u}\,\BC(P,q_u)\le\Bigl(\sum_u\sqrt{\lambda_u}\Bigr)\max_u\beta(P;A_u|B_u),
\]
and $\sum_u\sqrt{\lambda_u}\le\sqrt k$ by the Cauchy--Schwarz inequality.
\end{proof}

Mass that $Q$ puts outside the support of $P$ is allowed and never helps. Different terms may use different splits, because the final bound is uniform over balanced splits. This is why determinism is not needed.

\begin{lemma}[balanced decomposition~{\cite[Lemma~7.2]{Blaschke26Dimers}}]\label{lem:decomp}
Let $C$ be a smooth decomposable PC with nonnegative weights over $x\in\{0,1\}^E$, $\abs E=M\ge2$, whose leaves have arity at most $2M/3$. Then, as functions on $\{0,1\}^E$,
\[
  f_C=\sum_{u\in\cU}f_u(x_{A_u})\,h_u(x_{E\setminus A_u}),
\]
where each $f_u,h_u\ge0$, each $A_u$ satisfies $M/3<\abs{A_u}\le2M/3$, and $\abs\cU\le\abs C$, with size counted as nodes plus wires.
\end{lemma}

\begin{proof}
We repeat the proof of~\cite{Blaschke26Dimers} with $M$ in place of $n$. Let $X$ be the scope of the root. If $\abs X\le2M/3$, pick any $A\supseteq X$ with $M/3<\abs A\le2M/3$; then $f_C=f_C(x_A)\cdot1$ is one term. Otherwise, replace each product node of fan-in $d$ by a chain of $d-1$ binary products. This preserves decomposability and the computed functions, and the number of nodes becomes at most the original number of nodes plus wires, that is, at most $\abs C$. Let $\cU$ be the set of nodes $u$ with $\abs{X_u}\in(M/3,2M/3]$, and put $A_u:=X_u$. We claim that every node $v$ with $\abs{X_v}>2M/3$ satisfies
\[
  f_v=\sum_{u\in\cU}f_u\,h_{u,v},\qquad h_{u,v}\ge0\ \text{a function of }x_{X_v\setminus X_u},
\]
and argue by induction in topological order. No leaf has scope larger than $2M/3$. At a sum node, smoothness gives $X_c=X_v$ for every child $c$, so $f_v=\sum_cw_cf_c=\sum_uf_u\bigl(\sum_cw_ch_{u,c}\bigr)$. At a product node $v=c_1\times c_2$ with $\abs{X_{c_1}}\ge\abs{X_{c_2}}$ we have $\abs{X_{c_1}}>M/3$. If $\abs{X_{c_1}}\le2M/3$, then $c_1\in\cU$ and $f_v=f_{c_1}f_{c_2}$. Otherwise $f_v=\sum_uf_u\,(h_{u,c_1}f_{c_2})$, and $(X_{c_1}\setminus X_u)\sqcup X_{c_2}=X_v\setminus X_u$ by decomposability. Apply the claim at the root. A function of $x_{X\setminus X_u}$ is also a function of $x_{E\setminus X_u}$.
\end{proof}

This is the standard balancing argument for syntactically multilinear circuits~\cite{RazYehudayoff08}. Martens and Medabalimi state the same form with at most $s^2$ terms for decomposable sum-product networks of size $s$ with univariate leaves~\cite[Thm.~38]{MartensMedabalimi14}; using their version would halve the exponent in our bounds.

\begin{remark}[counting nodes only]\label{rem:nodes}
The prefactor $\abs\cU\le\abs C$ uses the convention that size counts nodes plus wires, because binarizing a product node of fan-in $d$ creates $d-2$ new nodes. If size counts nodes only, the same bound $\abs\cU\le\#\text{nodes}$ holds without binarizing. At a product node $v$ with $\abs{X_v}>2M/3$, at most one child is high (scope above $2M/3$) and a high child excludes a medium one, since scopes are disjoint and add up to at most $M$. If a child is high, expand it as above. If a child is medium, select it. If all children are low (scope at most $M/3$), add children to a group one at a time until the union of their scopes first exceeds $M/3$; the union then has size at most $2M/3$. The product of the group is a nonnegative function of a medium scope, and we give it the label $v$. Every label is then an original node, so there are at most as many terms as nodes.
\end{remark}

\begin{corollary}[circuit reduction]\label{cor:reduction}
Let $P$ be a distribution on $\{0,1\}^E$ and let $C$ be as in Lemma~\ref{lem:decomp}. If $1-\TV(P,Q_C)\ge\delta>0$, then
\[
  \abs C\ \ge\ \delta^2\Big/\max_{A|B\ \mathrm{balanced}}\beta(P;A|B)^2 .
\]
\end{corollary}

\begin{proof}
Lemma~\ref{lem:decomp}, applied to $f_C/\sum_xf_C(x)$, writes $Q_C$ as a sum of $k\le\abs C$ nonnegative products across balanced splits ($\abs{A_u}>M/3$ and $\abs{E\setminus A_u}\ge M/3$). Lemma~\ref{lem:overlap} gives $\delta\le\sqrt k\max\beta$.
\end{proof}

\section{Cycle bits and the local lemma}\label{sec:trees}

This section proves the bound on $\beta(\UST(G);A|B)$ for a single split, in terms of a family of crossing cycles. It is the only part of the paper that uses spanning trees.

Fix a split $E=A\sqcup B$. A cycle $C$ is \emph{mixed} if $E(C)$ meets both $A$ and $B$. For a mixed cycle define the two bits
\[
  u_C(x_A):=\one\bigl[E(C)\cap A\subseteq x_A\bigr],\qquad v_C(x_B):=\one\bigl[E(C)\cap B\subseteq x_B\bigr],
\]
where we identify $x_A$ with the set of edges of $A$ it contains. The bit $u_C$ depends only on the $A$-part of $x$ and $v_C$ only on the $B$-part. If $x$ is a spanning tree, then $u_C(x)=v_C(x)=1$ would mean $E(C)\subseteq x$, which is impossible. So under $\UST(G)$ the cell $(1,1)$ has probability zero.

\subsection{The spatial Markov property at a cycle}

Recall that $I(C)$ is the set of edges with an endpoint in $V(C)$. The next lemma describes the conditional law of $T\cap I(C)$ given the rest of the tree. It is a standard consequence of the uniformity of $T$; we include the proof.

\begin{lemma}[conditional law near a cycle]\label{lem:contract}
Let $C$ be a cycle of $G$ and let $F\subseteq E\setminus I(C)$ satisfy $\Pr[T\setminus I(C)=F]>0$. Then $F$ is a forest in which every vertex of $C$ is isolated. Let $K_1,\dots,K_s$ be the vertex sets of the connected components of the graph $(V\setminus V(C),F)$. Let $\mathcal M_F$ be the multigraph with node set $V(C)\cup\{K_1,\dots,K_s\}$ and one edge for each $e\in I(C)$, joining the nodes that contain the endpoints of $e$. Then $\mathcal M_F$ has no loops, every edge of $\mathcal M_F$ has an endpoint in $V(C)$, and conditionally on $T\setminus I(C)=F$, the set $T\cap I(C)$ is uniformly distributed over the edge sets of spanning trees of $\mathcal M_F$.
\end{lemma}

\begin{proof}
$F$ is contained in a tree, so it is a forest. Every edge at a vertex of $C$ lies in $I(C)$, so these vertices are isolated in $F$, and every edge of $I(C)$ has an endpoint in $V(C)$. Such an endpoint forms a node of its own, so no edge of $\mathcal M_F$ is a loop.

Since $T$ is uniform and $T=F\cup J$ with $J:=T\cap I(C)$, the conditional law of $J$ is uniform over the sets $J\subseteq I(C)$ for which $F\cup J$ is a spanning tree. We show that these are exactly the spanning trees of $\mathcal M_F$. The graph $(V,F)$ has $\abs{V(C)}+s$ components, so $\abs F=N-\abs{V(C)}-s$, and $\mathcal M_F$ has $\abs{V(C)}+s$ nodes. Each component of $F$ is connected, so $F\cup J$ is connected if and only if $J$ connects $\mathcal M_F$. A connected graph on $N$ vertices is a tree if and only if it has $N-1$ edges, and $\abs F+\abs J=N-1$ if and only if $\abs J=\abs{V(C)}+s-1$. So $F\cup J$ is a spanning tree if and only if $J$ is a connected spanning subgraph of $\mathcal M_F$ with one edge fewer than the number of nodes, that is, a spanning tree of $\mathcal M_F$.
\end{proof}

Let $d_j\ge0$ be the number of edges of $\mathcal M_F$ at the node $K_j$. Since $\mathcal M_F$ has a spanning tree, $d_j\ge1$ for every $j$.

\begin{lemma}[local completion]\label{lem:local}
Let $C$ be a cycle of $G$, let $e\in E(C)$, and let $F$ be as in Lemma~\ref{lem:contract}. For every $x_0\in V(C)$,
\[
  \Pr\bigl[\,T\cap E(C)=E(C)\setminus\{e\}\ \bigm|\ T\setminus I(C)=F\,\bigr]\ \ge\ \prod_{y\in V(C)\setminus\{x_0\}}\frac1{\deg(y)}\ \ge\ D^{-(\abs C-1)} .
\]
\end{lemma}

\begin{proof}
By Lemma~\ref{lem:contract} the conditional probability is the fraction of spanning trees $J$ of $\mathcal M_F$ with $J\cap E(C)=E(C)\setminus\{e\}$.

\emph{Numerator.} Let $J$ consist of the path $E(C)\setminus\{e\}$ together with exactly one edge of $\mathcal M_F$ at each node $K_j$. The path spans $V(C)$. Every edge at $K_j$ joins $K_j$ to a vertex of $V(C)$, so $J$ is connected, and it has $\abs C-1+s$ edges, one fewer than the number of nodes. So $J$ is a spanning tree. It contains no chord of $C$, so $J\cap E(C)=E(C)\setminus\{e\}$. Distinct choices give distinct sets, so there are $\prod_jd_j$ such trees.

\emph{Denominator.} Root a spanning tree $J$ of $\mathcal M_F$ at $x_0$. Every node other than $x_0$ has a unique parent edge in $J$, which is an edge of $\mathcal M_F$ at that node, and $J$ is the set of parent edges. So the map from $J$ to its list of parent edges is injective. The node $K_j$ has $d_j$ possible parent edges. A vertex $y\in V(C)\setminus\{x_0\}$ has at most $\deg(y)$, because all edges of $G$ at $y$ lie in $I(C)$ and these are its edges in $\mathcal M_F$. So $\mathcal M_F$ has at most $\prod_jd_j\prod_{y\ne x_0}\deg(y)$ spanning trees.

Dividing gives the first inequality, and $\abs{V(C)\setminus\{x_0\}}=\abs C-1$ gives the second.
\end{proof}

Chords of $C$ only enlarge the denominator's choice sets, and parallel edges of $\mathcal M_F$ are kept distinct throughout. The case $s=0$ uses empty products. The bound is uniform over all outside configurations $F$, and this uniformity is what the chain rule needs.

\begin{remark}
The lemma conditions on $T$ outside \emph{all} edges incident to $C$, not just outside $E(C)$. This matters. If one conditioned only on $T\setminus E(C)$, the cycle could freeze. For example, let $C$ be a triangle and let a further vertex $w$ be joined to all three vertices of $C$. Given that all three edges at $w$ are in $T$, no edge of $C$ can be. Lemma~\ref{lem:local} leaves those three edges free.
\end{remark}

\subsection{The bound for one split}

\begin{theorem}[affinity bound for a split]\label{thm:cut}
Let $G$ be a connected graph with maximum degree $D$, let $E=A\sqcup B$ be any split of its edges, and let $C_1,\dots,C_r$ be pairwise vertex-disjoint mixed cycles. For each $k$ put $p_k:=\max_{x_0\in V(C_k)}\prod_{y\in V(C_k)\setminus\{x_0\}}\deg(y)^{-1}$. Then
\[
  \beta\bigl(\UST(G);A|B\bigr)\ \le\ \prod_{k=1}^r\Bigl(\frac{1+\sqrt{1-4p_k^2}}2\Bigr)^{1/2}\ \le\ \prod_{k=1}^r\bigl(1-p_k^2\bigr)^{1/2}.
\]
In particular, if every $C_k$ has length at most $\ell$, then $\beta(\UST(G);A|B)\le\bigl(1-D^{-2(\ell-1)}\bigr)^{r/2}$.
\end{theorem}

\begin{proof}
Let $T\sim\UST(G)$ and put $u_k:=u_{C_k}(T\cap A)$ and $v_k:=v_{C_k}(T\cap B)$. The vector $u=(u_1,\dots,u_r)$ is a function of $T\cap A$ and $v$ is a function of $T\cap B$. By Lemma~\ref{lem:dpi}, $\beta(\UST(G);A|B)\le\opnorm{\sqrt\mu}$, where $\mu$ is the joint law of $(u,v)$ with rows indexed by $u$ and columns by $v$. Order the pairs as $(u_1,v_1),\dots,(u_r,v_r)$; the order is fixed before $T$ is sampled. By Lemmas~\ref{lem:chain} and~\ref{lem:two-by-two} it suffices to show that for every $k$ and every value $\pi$ of the past $(u_{<k},v_{<k})$ with positive probability,
\[
  \Pr\bigl[(u_k,v_k)=(1,1)\bigm|\pi\bigr]=0,\qquad\Pr\bigl[(u_k,v_k)=(0,1)\bigm|\pi\bigr]\ge p_k,\qquad\Pr\bigl[(u_k,v_k)=(1,0)\bigm|\pi\bigr]\ge p_k .
\]
The first holds because $T$ contains no cycle. For the others, let $F_k:=T\setminus I(C_k)$. An edge of $C_j$ with $j\ne k$ has both endpoints in $V(C_j)$, which is disjoint from $V(C_k)$, so $E(C_j)\cap I(C_k)=\emptyset$. Hence the past $(u_{<k},v_{<k})$ is a function of $F_k$, and the event $\{(u_{<k},v_{<k})=\pi\}$ is a union of events $\{F_k=F\}$. Pick $e_A\in E(C_k)\cap A$ and $e_B\in E(C_k)\cap B$. If $T\cap E(C_k)=E(C_k)\setminus\{e_A\}$, then $e_A\notin T$ while every edge of $E(C_k)\cap B$ is in $T$, so $(u_k,v_k)=(0,1)$. Likewise $E(C_k)\setminus\{e_B\}$ gives $(1,0)$. By Lemma~\ref{lem:local}, both events have conditional probability at least $p_k$ given every value $F$ of $F_k$ with positive probability. By the tower property the same bounds hold given $\pi$, since conditioning on $\pi$ averages the conditional laws given $F_k=F$ over the values $F$ compatible with $\pi$.
\end{proof}

The incident sets $I(C_j)$ and $I(C_k)$ may share an edge joining the two cycles. This does not matter, because the bits use cycle edges only. The theorem holds for every split, balanced or not, and every family of vertex-disjoint mixed cycles. The geometric input of Section~\ref{sec:geometry} is a lower bound on $r$ for balanced splits.

\begin{corollary}[tree rectangles]\label{cor:rectangle}
Let $R_A\subseteq2^A$ and $R_B\subseteq2^B$ be nonempty, and suppose that every set $a\cup b$ with $a\in R_A$, $b\in R_B$ is a spanning tree. Let $R$ be the set of these trees. Then $\UST(G)(R)\le\beta(\UST(G);A|B)^2$.
\end{corollary}

\begin{proof}
The map $(a,b)\mapsto a\cup b$ is injective, since $A\cap B=\emptyset$. Let $q$ be the product of the uniform distributions on $R_A$ and $R_B$, so $q$ is uniform on $R$. Every element of $R$ has probability $1/\tau(G)$ under $P=\UST(G)$, so $\BC(P,q)=\abs R/\sqrt{\tau(G)\abs R}=\sqrt{P(R)}$. Now use $\BC(P,q)\le\beta$.
\end{proof}

\section{Geometry: many disjoint crossing cycles}\label{sec:geometry}

Theorem~\ref{thm:cut} turns vertex-disjoint mixed short cycles into a bound on $\beta$. This section shows that, in the graphs of Theorem~\ref{thm:main}, every balanced split has linearly many of them. The condition we need is that short cycles tie together the edges at every vertex.

\begin{definition}[linked graphs]\label{def:linked}
Let $G$ be a connected graph with $N\ge3$ vertices, and let $\cC$ be a nonempty family of cycles of $G$. For a vertex $x$, let $\Lambda_x$ be the graph whose nodes are the edges at $x$, with $e\sim f$ when some cycle of $\cC$ contains both $e$ and $f$. The family $\cC$ \emph{links} $G$ if $\Lambda_x$ is connected for every vertex $x$. We write $\ell:=\max_{C\in\cC}\abs C$ and $\kappa:=\max_x\abs{\{C\in\cC:x\in V(C)\}}$.
\end{definition}

\begin{lemma}\label{lem:linked-basic}
If $\cC$ links $G$, then every edge of $G$ lies on a cycle of $\cC$. In particular $G$ has no bridge, its minimum degree is at least $2$, and $1\le\kappa<\infty$.
\end{lemma}

\begin{proof}
Let $e=xy$. Since $G$ is connected with at least three vertices, $x$ or $y$ has degree at least $2$, say $x$. Then $\Lambda_x$ is connected with at least two nodes, so $e$ is adjacent in $\Lambda_x$ to some other edge, which means that some cycle of $\cC$ contains $e$. An edge on a cycle is not a bridge, and every vertex has an edge on a cycle, hence degree at least $2$.
\end{proof}

The conditions $N\ge3$ and $\cC\ne\emptyset$ exclude $K_2$, which is vacuously linked by the empty family and has a bridge. Some condition of this kind is needed: Proposition~\ref{prop:pendant} shows that pendant edges defeat any bound that is uniform over balanced splits.

\subsection{Mixed cycles in linked expanders}

\begin{lemma}[mixed cycles]\label{lem:mixed}
Let $G$ be a connected graph with $N\ge3$ vertices, maximum degree $D$ and edge expansion $h>0$, linked by a family $\cC$ with parameters $\ell,\kappa$. Colour the edges with two colours so that each colour class has at least $m_0>0$ edges, and let $t$ be the number of cycles of $\cC$ that contain edges of both colours. Then
\[
  t\ \ge\ \min\Bigl\{\frac{hN}{4D\ell},\ \frac{2m_0}{\ell D(1+D/h)}\Bigr\},
\]
and some $\lceil t/(\ell\kappa)\rceil$ of these mixed cycles are pairwise vertex-disjoint.
\end{lemma}

\begin{proof}
Let $Z$ be the set of vertices of mixed cycles of $\cC$. Then $\abs Z\le\ell t$.

\emph{Good vertices are monochromatic.} Let $x\notin Z$. Every cycle of $\cC$ through $x$ is monochromatic, so adjacent nodes of $\Lambda_x$ have the same colour. Since $\Lambda_x$ is connected, all edges at $x$ have one colour $c(x)$. If $x,y\notin Z$ are adjacent, then $c(x)=c(xy)=c(y)$. Hence for each connected component $K$ of $G-Z$, all edges with an endpoint in $K$ have the same colour.

\emph{Case $\abs Z\ge hN/(4D)$.} Then $t\ge\abs Z/\ell\ge hN/(4D\ell)$.

\emph{Case $\abs Z<hN/(4D)$.} Since $h\le D$, $\abs Z<N/4$. Suppose every component of $G-Z$ had at most $N/2$ vertices. If one of them has at least $N/4$ vertices, let $U$ be that component. Otherwise add components one at a time until the union first has at least $N/4$ vertices; this happens because $\abs{V\setminus Z}>3N/4$, and the union then has fewer than $N/2$ vertices. In both cases $U$ is a union of components with $N/4\le\abs U\le N/2$. Every edge of $\partial U$ has its other endpoint in $Z$, so
\[
  \frac{hN}4\le h\abs U\le\abs{\partial U}\le D\abs Z<\frac{hN}4,
\]
a contradiction. Hence $G-Z$ has a component $K$ with $\abs K>N/2$. Let $S:=V\setminus(K\cup Z)$. Then $\abs S<N/2$, and every edge of $\partial S$ goes to $Z$, so $h\abs S\le D\abs Z$. Therefore
\[
  \abs{V\setminus K}=\abs S+\abs Z\le(1+D/h)\abs Z .
\]
All edges with an endpoint in $K$ have the colour of $K$, so every edge of the other colour lies inside $V\setminus K$. There are at most $D\abs{V\setminus K}/2\le\ell D(1+D/h)t/2$ such edges, and at least $m_0$. This gives the second bound.

\emph{Packing.} Choose mixed cycles greedily. A chosen cycle shares a vertex with at most $\ell\kappa$ cycles of $\cC$, itself included, because it has at most $\ell$ vertices and each lies on at most $\kappa$ cycles. Discard these and repeat. This produces at least $\lceil t/(\ell\kappa)\rceil$ pairwise vertex-disjoint mixed cycles.
\end{proof}

\subsection{Prisms and squares}

\begin{lemma}[prisms]\label{lem:prism}
Let $H$ be a connected simple graph with $n\ge3$ vertices, maximum degree at most $\Delta$ and edge expansion $h(H)\ge h_0>0$, and let $G=H\square K_2$. Then $N=2n$, $M=2\abs{E(H)}+n\ge N$, $D\le\Delta+1$, and $h(G)\ge\min(1,h_0)/3$. The family of $4$-cycles
\[
  Q_{uv}:=(u,0)\,(v,0)\,(v,1)\,(u,1),\qquad uv\in E(H),
\]
links $G$ with $\ell=4$ and $\kappa\le\Delta$.
\end{lemma}

\begin{proof}
The counts are immediate, and $M\ge2(n-1)+n\ge2n$ since $H$ is connected and $n\ge2$.

\emph{Linking.} The edges at $(v,0)$ are the rung $(v,0)(v,1)$ and the edges $(v,0)(u,0)$ for $u\sim v$ in $H$. Each $(v,0)(u,0)$ lies in $Q_{uv}$ together with the rung. So $\Lambda_{(v,0)}$ contains a spanning star centred at the rung, and is connected. The same holds at $(v,1)$. The vertex $(v,i)$ lies on $Q_{uv}$ for the $\deg_H(v)\le\Delta$ neighbours $u$ of $v$, and on no other cycle of the family.

\emph{Expansion.} Let $S=S_0\times\{0\}\cup S_1\times\{1\}$ with $0<\abs S\le n$, and put $x=\abs{S_0}$, $y=\abs{S_1}$. By symmetry of the two layers we may assume $x\ge y$. If $x\le n/2$, the edges of the two copies of $H$ give $\abs{\partial S}\ge h_0x+h_0y$. If $x>n/2$, then $y\le n-x<n/2$, so the copy of $H$ in layer $1$ contributes at least $h_0y$ boundary edges, and each of the at least $x-y$ vertices $v\in S_0\setminus S_1$ contributes its rung. So $\abs{\partial S}\ge h_0y+(x-y)$. If $2y\le x$, this is at least $x/2\ge(x+y)/3$. If $2y>x$, it is at least $h_0y\ge h_0(x+y)/3$. In all cases $\abs{\partial S}\ge\min(1,h_0)\abs S/3$.
\end{proof}

\begin{lemma}[squares]\label{lem:square}
Let $H$ be as in Lemma~\ref{lem:prism} and let $G=H^2$. Then $N=n$, $M\ge N-1\ge2N/3$, $D\le\Delta^2$, and $h(G)\ge h_0$. The family of triangles $\{u,v,w\}$ such that $u\,v\,w$ is a path in $H$ links $G$ with $\ell=3$ and $\kappa\le3\Delta(\Delta-1)/2$.
\end{lemma}

\begin{proof}
A vertex has at most $\Delta$ neighbours in $H$ and at most $\Delta(\Delta-1)$ vertices at distance two, so $D\le\Delta^2$. $G$ contains $H$ on the same vertex set, so $\abs{\partial_GS}\ge\abs{\partial_HS}$ for every $S$, and $h(G)\ge h(H)\ge h_0$. $G$ is connected, so $M\ge N-1$, and $N-1\ge2N/3$ because $N\ge3$.

\emph{Linking.} For a path $u\,v\,w$ in $H$, the vertices $u,v,w$ are pairwise at distance at most two, so they span a triangle of $G$. Let $x$ be a vertex. Any two $H$-edges $xu$, $xw$ at $x$ lie in the triangle of the path $u\,x\,w$. An edge $xw$ of $G$ that is not in $H$ joins vertices at distance two, so there is a path $x\,v\,w$ in $H$, and its triangle contains $xw$ and the $H$-edge $xv$. So every node of $\Lambda_x$ is adjacent to an $H$-edge at $x$, the $H$-edges at $x$ are pairwise adjacent, and $\Lambda_x$ is connected. The family is nonempty because $H$ is connected with $n\ge3$.

\emph{The parameter $\kappa$.} Assign to each triangle of the family one path that generates it. A triangle through $x$ is assigned a path with centre $x$ or a path with endpoint $x$. There are at most $\binom\Delta2$ paths of the first kind and at most $\Delta(\Delta-1)$ of the second, so $\kappa\le3\Delta(\Delta-1)/2$.
\end{proof}

\subsection{The square grid}

The grid is not an expander, but it is linked by its unit squares, and an isoperimetric inequality can replace expansion. Let $G_L$ be the $L\times L$ grid with $L\ge2$, so $N=L^2$, $M=2L(L-1)$ and $D=4$. At an interior vertex, the four edges and the four unit squares around it form a $4$-cycle in $\Lambda_x$; at boundary vertices $\Lambda_x$ is a path. So the unit squares link $G_L$ with $\ell=4$ and $\kappa=4$.

\begin{lemma}[grid isoperimetry]\label{lem:grid-iso}
For every $S\subseteq V(G_L)$ with $0<\abs S\le L^2/2$, $\abs{\partial S}\ge\min\{2\sqrt{\abs S},\,L/2\}$.
\end{lemma}

\begin{proof}
Call a row full if it lies in $S$, empty if it misses $S$, and partial otherwise; the same for columns. A partial row contains two horizontally adjacent vertices, one in $S$ and one not, so it contributes a horizontal edge to $\partial S$; a partial column contributes a vertical one. Let $a$ and $b$ be the numbers of partial rows and columns.
If some row is full and some row is empty, every column is partial, so $\abs{\partial S}\ge L$. If some row is full and no row is empty, then at most $\abs S/L\le L/2$ rows are full, so $a\ge L/2$. The same two cases apply to columns. In the remaining case no row and no column is full, so the rows and columns that meet $S$ are partial, $\abs S\le ab\le((a+b)/2)^2$, and $\abs{\partial S}\ge a+b\ge2\sqrt{\abs S}$.
\end{proof}

\begin{lemma}[mixed squares in the grid]\label{lem:grid-mixed}
Let $L\ge2$. Colour the edges of $G_L$ with two colours so that each colour class has at least $M/3$ edges. Then at least $L/32$ unit squares contain edges of both colours, and some $\lceil L/512\rceil$ of them are pairwise vertex-disjoint.
\end{lemma}

\begin{proof}
Let $Z$ be the set of vertices of mixed unit squares and $t$ their number, so $\abs Z\le4t$. As in Lemma~\ref{lem:mixed}, each component of $G_L-Z$ has all its incident edges of one colour. We show $\abs Z\ge L/8$, which gives $t\ge L/32$. Suppose $\abs Z<L/8$. Then $\abs Z<N/4$. If every component of $G_L-Z$ had at most $N/2$ vertices, the argument of Lemma~\ref{lem:mixed} gives a union $U$ of components with $N/4\le\abs U\le N/2$ and $\abs{\partial U}\le4\abs Z<L/2$. Lemma~\ref{lem:grid-iso} gives $\abs{\partial U}\ge\min\{2\sqrt{N/4},L/2\}=L/2$, a contradiction. So there is a component $K$ with $\abs K>N/2$. Let $U_1,U_2,\dots$ be the other components. Each has fewer than $N/2$ vertices, the sets $\partial U_i$ are disjoint, and all their edges end in $Z$, so $\sum_i\abs{\partial U_i}\le4\abs Z<L/2$. By Lemma~\ref{lem:grid-iso}, $\abs{U_i}\le\abs{\partial U_i}^2/4$, hence $\sum_i\abs{U_i}\le(\sum_i\abs{\partial U_i})^2/4\le4\abs Z^2$. Every edge of the colour not used at $K$ lies inside $V\setminus K$, which has at most $\abs Z+4\abs Z^2$ vertices, so there are at most $2\abs Z+8\abs Z^2<L/4+L^2/8$ such edges. But there are at least $M/3=2L(L-1)/3$, and $L/4+L^2/8\le2L(L-1)/3$ for $L\ge2$. This is a contradiction. The packing step is as in Lemma~\ref{lem:mixed} with $\ell\kappa=16$.
\end{proof}

\section{Main results}\label{sec:main}

\subsection{Proof of Theorem~\ref{thm:main} and the families it covers}

\begin{proof}[Proof of Theorem~\ref{thm:main}]
By Theorem~\ref{thm:touch}, $\beta(P;A|B)^2\le e^{-cN}$ for every balanced split. Apply Corollary~\ref{cor:reduction} with $\delta=1-\eps>0$.
\end{proof}

Exactly as for Theorem~\ref{thm:general} below, the bound stays exponential when the overlap $\delta=1-\TV$ is itself exponentially small, as long as $\delta^2\ge e^{-(c-c')N}$ for some $c'>0$.

Below, \emph{the easy Cheeger direction} is Lemma~\ref{lem:easy-cheeger}: $h(G)\ge\lambda_2(L_G)/2$.

\begin{corollary}[families]\label{cor:families}
Each of the following families satisfies the hypotheses of Theorem~\ref{thm:main} with $D$ and $h$ independent of $N$.
\begin{enumerate}[label=(\alph*)]
\item For fixed $d\ge3$, a uniformly random $d$-regular graph on $n$ vertices, with probability tending to one as $n\to\infty$ ($nd$ even).
\item The Ramanujan graphs $X^{p,q}$ of Lubotzky, Phillips and Sarnak~\cite{LPS88}, with $D=p+1$ and $h=(\sqrt p-1)^2/2$, for $q$ large enough that they are simple graphs (their girth tends to infinity with $q$~\cite{LPS88}).
\item The prisms $H\square K_2$ and squares $H^2$ of connected graphs $H$ with at least three vertices, maximum degree at most $\Delta$ and $h(H)\ge h_0>0$, with $D$ and $h$ as in Lemmas~\ref{lem:prism} and~\ref{lem:square}.
\item Every 2-connected graph $H$ with maximum degree at most $\Delta$ and $h(H)\ge h_0>0$, itself.
\end{enumerate}
\end{corollary}

\begin{proof}
(a) For fixed $d\ge3$, a random $d$-regular graph is $d$-connected with probability tending to one~\cite{Wormald81}, and there is $h_d>0$ such that its edge expansion is at least $h_d$ with probability tending to one~\cite{Bollobas88}.

(b) The graphs $X^{p,q}$ are connected $(p+1)$-regular Cayley graphs, and every adjacency eigenvalue other than $\pm(p+1)$ has absolute value at most $2\sqrt p$~\cite{LPS88}. So $\lambda_2(L)\ge p+1-2\sqrt p=(\sqrt p-1)^2$, and $h\ge(\sqrt p-1)^2/2$ by the easy Cheeger direction. A Cayley graph is vertex-transitive. Every connected graph with at least two vertices has a vertex that is not a cut vertex, for instance a leaf of a spanning tree, and in a vertex-transitive graph every vertex is then a non-cut vertex. So a connected vertex-transitive graph with at least three vertices is 2-connected.

(c) The expansion and degree bounds are Lemmas~\ref{lem:prism} and~\ref{lem:square}. For 2-connectivity, delete a vertex $(x,0)$ from $H\square K_2$: the copy $H\times\{1\}$ is still connected, and every remaining $(y,0)$ is joined to $(y,1)$; the case $(x,1)$ is symmetric. Delete a vertex $v$ from $H^2$: every component of $H-\{v\}$ contains an $H$-neighbour of $v$, since $H$ is connected, and the $H$-neighbours of $v$ are pairwise adjacent in $H^2$, so $H^2-\{v\}$ is connected.

(d) There is nothing to prove.
\end{proof}

For the graphs of (b) with $p\ge5$, $\lambda_2(L)\ge(\sqrt p-1)^2>1$, so Theorem~\ref{thm:main} holds with the spectral constant $c_\lambda(p+1)$, $\lambda=(\sqrt p-1)^2$; for $p=5$ this is $8.0\cdot10^{-8}$. The same applies to random $d$-regular graphs with $d\ge6$, with $c_\lambda(d)$ for any fixed $\lambda\in(1,d-2\sqrt{d-1})$, by Friedman's theorem~\cite{Friedman08} that all adjacency eigenvalues other than $d$ have absolute value at most $2\sqrt{d-1}+o(1)$ with probability tending to one.

The Ramanujan graphs of~\cite{LPS88} have girth of order $\log N$, and random regular graphs have few short cycles: for fixed $d$ and $k$, the number of cycles of length $k$ converges in distribution to a Poisson variable with mean $(d-1)^k/(2k)$~\cite{Wormald81cycles}. The local argument of Sections~\ref{sec:trees} and~\ref{sec:geometry} gives nothing there (Section~\ref{sec:limits}), and an earlier version of this paper left this case open.

\subsection{Linked expanders}

\begin{theorem}[linked expanders]\label{thm:general}
Let $G$ be a connected simple graph with $N\ge3$ vertices, $M$ edges, maximum degree $D$ and edge expansion $h>0$, linked by a family of cycles with parameters $\ell$ and $\kappa$. Put
\[
  r_0:=\frac1{\ell\kappa}\min\Bigl\{\frac{hN}{4D\ell},\ \frac{2M}{3\ell D(1+D/h)}\Bigr\}.
\]
Then $\beta(\UST(G);A|B)^2\le(1-D^{-\ell})^{r_0}$ for every balanced split $A|B$. Consequently, if $C$ is a smooth decomposable PC with nonnegative weights over the $M$ edge variables, with leaves of arity at most $2M/3$, and $1-\TV(\UST(G),Q_C)\ge\delta>0$, then
\[
  \abs C\ \ge\ \delta^2(1-D^{-\ell})^{-r_0}\ \ge\ \delta^2\exp\bigl(r_0\,D^{-\ell}\bigr).
\]
\end{theorem}

\begin{proof}
Let $A|B$ be balanced and colour $A$ and $B$ with two colours. Each class has at least $M/3$ edges, so Lemma~\ref{lem:mixed} with $m_0=M/3$ gives at least $r_0$ pairwise vertex-disjoint mixed cycles of length at most $\ell$. Theorem~\ref{thm:cut} gives $\beta^2\le(1-D^{-\ell})^{r_0}$. Corollary~\ref{cor:reduction} gives the circuit bound, and $-\log(1-x)\ge x$ gives the last inequality.
\end{proof}

With $\delta=1-\eps$ and $c:=r_0D^{-\ell}/N$ this is the bound $\abs C\ge(1-\eps)^2e^{cN}$. The bound stays exponential even when the overlap $\delta$ is itself exponentially small, as long as $\delta^2\ge e^{-(c-c')N}$ for some $c'>0$.

\subsection{Prisms and squares: proof of Corollary~\ref{cor:prism-square}}

\begin{corollary}[closed-form constants]\label{cor:constants}
Let $H$ satisfy the hypotheses of Corollary~\ref{cor:prism-square} and put $h'=\min(1,h_0)$. For $G=H\square K_2$,
\[
  r_0\,D^{-\ell}\ \ge\ \frac{h'}{24\,\Delta(\Delta+1)^5(3\Delta+4)}\,N ,
\]
and for $G=H^2$,
\[
  r_0\,D^{-\ell}\ \ge\ \frac{8h_0}{243\,\Delta^{9}(\Delta-1)(h_0+\Delta^2)}\,N .
\]
\end{corollary}

\begin{proof}
The quantity $r_0D^{-\ell}$ increases with $h$ and $M$ and decreases with $D$ and $\kappa$, so we may insert the bounds of Lemmas~\ref{lem:prism} and~\ref{lem:square}.

\emph{Prism.} Here $D\le\Delta+1$, $\ell=4$, $\kappa\le\Delta$, $h\ge h'/3$, $M\ge N$ and $D^{-\ell}\ge(\Delta+1)^{-4}$. The first term of the minimum is at least $h'N/(48(\Delta+1))$. The second is at least
\[
  \frac{2N}{12(\Delta+1)(1+3(\Delta+1)/h')}=\frac{h'N}{6(\Delta+1)(h'+3\Delta+3)}\ge\frac{h'N}{6(\Delta+1)(3\Delta+4)},
\]
using $h'\le1$. Since $6(3\Delta+4)\ge48$ for $\Delta\ge2$, the minimum is at least the last expression. Dividing by $\ell\kappa\le4\Delta$ and multiplying by $D^{-\ell}$ gives the claim.

\emph{Square.} Here $D\le\Delta^2$, $\ell=3$, $\kappa\le3\Delta(\Delta-1)/2$, $h\ge h_0$, $M\ge2N/3$ and $D^{-\ell}\ge\Delta^{-6}$. The first term is at least $h_0N/(12\Delta^2)$, and the second is at least
\[
  \frac{4N}{27\Delta^2(1+\Delta^2/h_0)}=\frac{4h_0N}{27\Delta^2(h_0+\Delta^2)} .
\]
The second is the smaller, since $27(h_0+\Delta^2)\ge48$. Dividing by $\ell\kappa\le9\Delta(\Delta-1)/2$ and multiplying by $D^{-\ell}$ gives the claim.
\end{proof}

\begin{proof}[Proof of Corollary~\ref{cor:prism-square}]
$G$ is 2-connected by Corollary~\ref{cor:families}(c). Lemmas~\ref{lem:prism} and~\ref{lem:square} show that $G$ satisfies the hypotheses of Theorem~\ref{thm:general}. Apply it with $\delta=1-\eps$ and use Corollary~\ref{cor:constants}.
\end{proof}

The closed forms are lossy. A direct evaluation of $r_0$ is better, as the next example shows.

\begin{example}[cubic prisms]\label{ex:cubic}
Let $H$ be $3$-regular with $h(H)\ge1/5$, and $G=H\square K_2$. Then $D=4$, $M=2N$, $\ell=4$, $\kappa=3$, $h\ge1/15$ and $D^{-\ell}=1/256$. The two terms in $r_0$ are $hN/(4D\ell)\ge N/960$ and $2M/(3\ell D(1+D/h))\ge N/732$, so $r_0\ge N/11520$. The per-cycle factor in Theorem~\ref{thm:cut} is at most $(1-4^{-4})^{1/2}=0.9980449\ldots$; the sharper form $((1+\sqrt{1-4^{-3}})/2)^{1/2}=0.9980372\ldots$ changes the fifth decimal. So
\[
  \abs C\ \ge\ (1-\eps)^2\,(1-4^{-4})^{-N/11520}\ \ge\ (1-\eps)^2\,e^{3.3974\cdot10^{-7}N}\ \ge\ (1-\eps)^2\,e^{3\cdot10^{-7}N}.
\]
Here $-\log(1-4^{-4})/11520=3.39748\ldots\cdot10^{-7}$. The closed form of Corollary~\ref{cor:constants} gives only $c\ge1/4792320\approx2.09\cdot10^{-7}$ for these parameters. Either way the bound exceeds $10$ only when $N$ exceeds about $7\cdot10^6$. The bound of the first version of this paper, with $D^{-2(\ell-1)}=4^{-6}$ in place of $4^{-4}$, was $16$ times weaker.
\end{example}

The constants are poor for three reasons: the per-cycle weight $D^{-\ell}$ is crude, the forbidden-cell factor uses two bits per cycle, and the geometric counting and packing lose factors of $\ell$, $D$ and $\kappa$. In exact computations on five small graphs, the geometric mean of the conditional factors $\gamma_k$ of Lemma~\ref{lem:chain} was between $0.94$ and $0.99$ per cycle, while the proved per-cycle bounds on the same instances were between $0.992$ and $0.998$ (Appendix~\ref{app:checks}). These graphs are far too small to say anything about rates. We made no attempt to optimize.

\subsection{The square grid}

\begin{corollary}[grid]\label{cor:grid}
Let $L\ge2$ and let $P$ be the uniform spanning tree of the $L\times L$ grid, with $M=2L(L-1)$ edges. Every smooth decomposable PC with nonnegative weights and leaves of arity at most $2M/3$ that satisfies $1-\TV(P,Q_C)\ge\delta>0$ has size at least $\delta^2(1-4^{-4})^{-L/512}\ge\delta^2e^{L/2^{17}}$.
\end{corollary}

\begin{proof}
Lemma~\ref{lem:grid-mixed} gives at least $L/512$ vertex-disjoint mixed unit squares for every balanced split. All degrees are at most $4$, so $q_k\ge4^{-4}$ in Theorem~\ref{thm:cut}. Corollary~\ref{cor:reduction} finishes the proof.
\end{proof}

The bound is tight up to the constant in the exponent. Process the edges in a fixed order, sweeping the grid with a staircase frontier of $O(L)$ active vertices. A state records the partition of the active vertices induced by the processed edges chosen so far, provided they form a forest in which every component still meets the frontier (a completed spanning component is allowed only at the end). Choosing an edge inside a block is rejected, choosing one across blocks merges them, and a component that would leave the frontier unconnected is rejected. Because the processed region is planar and lies on one side of the frontier, the realizable partitions are noncrossing, so there are at most a Catalan number $2^{O(L)}$ of them; recording arbitrary partitions would give only the weaker bound $2^{O(L\log L)}$. The resulting layered decision diagram has $O(L^2)\,2^{O(L)}$ nodes. Writing each node at layer $i$ as $[x_i=0]f_0+[x_i=1]f_1$, with all branches kept layered so both children have the same scope, gives a smooth, deterministic, decomposable circuit with univariate leaves of size $2^{O(L)}$ that computes the spanning-tree indicator; dividing by $\tau$ gives the uniform spanning tree. For the uniform domino tiling of the same grid, a bound $e^{\Theta(L)}$ was proved in~\cite{Blaschke26Dimers} by a much longer argument; Section~\ref{sec:limits} says why the short argument here does not apply to matchings.

Theorem~\ref{thm:main} does not cover the grid as stated: $h(G_L)$ is of order $1/L$, so $c(D,h)N$ tends to zero. Theorem~\ref{thm:touch-cut} itself needs no expansion, and we expect a local version to work. If $W$ has pairwise distances at least $3$, any two neighbours of $w\in W$ are joined in $G_L-W$ by a path of length at most $4$ around $w$, so $\rho_w\le4$ and $p_A(w)p_B(w)\ge6^{-4}$ by Lemma~\ref{lem:touch}. What remains is to show that every balanced split has $\Omega(L)$ mixed vertices, using Lemma~\ref{lem:grid-iso} in place of expansion. We have not written this out, so we state the vertex-touch bound for the grid only as a conjecture. It would not improve on Corollary~\ref{cor:grid} beyond constants.

\subsection{Upper bounds}

\begin{proposition}[upper bounds]\label{prop:upper}
Let $G$ be any connected graph with $N$ vertices and $M$ edges, and $P=\UST(G)$.
\begin{enumerate}[label=(\alph*)]
\item The generating polynomial $\sum_TP(T)\prod_{e\in T}z_e$ is computed by an arithmetic circuit with subtraction (a probabilistic generating circuit in the sense of~\cite{ZhangJubaVdB21}) of size polynomial in $N+M$.
\item For disjoint $S_0,S_1\subseteq E$, the probability $P(S_1\subseteq T,\ T\cap S_0=\emptyset)$ is a ratio of two determinants of integer matrices of order less than $N$, so it can be computed exactly in polynomial time.
\item Some smooth, deterministic, decomposable PC with univariate leaves represents $P$ exactly and has size $O(M\tau(G))=O(M2^M)$.
\end{enumerate}
\end{proposition}

\begin{proof}
(a) Let $L_z$ be the weighted Laplacian, with $L_z[x,x]=\sum_{e\ni x}z_e$ and $L_z[x,y]=-\sum_{e=xy}z_e$ for $x\ne y$, and let $L_z^{(v)}$ be $L_z$ with the row and column of a vertex $v$ removed. By the matrix-tree theorem~\cite{Kirchhoff1847}, $\det L_z^{(v)}=\sum_T\prod_{e\in T}z_e$. The algorithms of Berkowitz~\cite{Berkowitz84} and of Mahajan and Vinay~\cite{MahajanVinay99} compute the determinant of an $n\times n$ matrix by a division-free arithmetic circuit of size polynomial in $n$. Substituting the linear forms $L_z^{(v)}[x,y]$ and multiplying by the constant $1/\tau(G)$, where $\tau(G)=\det L_{\one}^{(v)}$, gives the circuit.

(b) If $S_1$ contains a cycle the probability is zero. Otherwise it equals $\tau(G')/\tau(G)$, where $G'$ is the multigraph obtained from $G$ by deleting $S_0$, contracting $S_1$ and removing loops. The matrix-tree theorem holds for multigraphs, so both counts are determinants of reduced Laplacians.

(c) Take a sum node with weight $1/\tau(G)$ on each spanning tree $T$, whose child is the product of the indicator leaves $[x_e=\one[e\in T]]$ over all $e\in E$.
\end{proof}

Since $N-1\le M\le DN/2$, parts (c) and Theorem~\ref{thm:main} together say that, for fixed $D$, $h$ and $\eps<1$, the optimal size of a nonnegative smooth decomposable PC that $\eps$-approximates $P$ is $2^{\Theta(M)}$. Part~(a) is not a nonnegative PC: the determinant circuit uses subtraction, and that is the point of the separation.

\subsection{Exact support}\label{sec:support}

Exact representation of a distribution with the same support as $P$ but different probabilities is not the case $\eps=0$ of Theorem~\ref{thm:main}, since such a distribution can be far from $P$ in total variation. Lower bounds for exact support follow instead from Corollary~\ref{cor:rectangle} and a rectangle cover.

\begin{corollary}[exact support]\label{cor:support}
Let $G$ satisfy the hypotheses of Theorem~\ref{thm:main}, and let $c$ and $\kappa=3c$ be as in Theorem~\ref{thm:touch}.
\begin{enumerate}[label=(\alph*)]
\item Let $C$ be a smooth decomposable PC with nonnegative weights and leaves of arity at most $2M/3$ such that $f_C(x)>0$ exactly when $x$ is a spanning tree. Then $\abs C\ge e^{cN}$.
\item Let $\Gamma$ be a monotone arithmetic circuit with fan-in at most two and $s$ nodes computing a polynomial whose monomials are exactly $\prod_{e\in T}z_e$, $T\in\cT(G)$, each with a positive coefficient. Then $s\ge e^{2cN/3}$.
\end{enumerate}
The same holds if $G$ satisfies the hypotheses of Theorem~\ref{thm:general} instead, with $c_0:=r_0D^{-\ell}/N$ in place of $c$ in (a), and in (b) with $c_1$ in place of $2c/3$, where $c_1$ is defined like $c_0$ with $M$ replaced by $N-1$ in the definition of $r_0$.
\end{corollary}

\begin{proof}
(a) Lemma~\ref{lem:decomp} writes $f_C=\sum_{u}f_uh_u$ with $k\le\abs C$ terms across balanced splits $A_u|B_u$. All terms are nonnegative, so the support of each term lies in $\cT(G)$, and the supports cover $\cT(G)$. The support of $f_uh_u$ is the rectangle $R_u=\{a\cup b:\ f_u(a)>0,\ h_u(b)>0\}$. By Corollary~\ref{cor:rectangle} and Theorem~\ref{thm:touch},
\[
  1\le\sum_u\UST(G)(R_u)\le k\max_u\beta(\UST(G);A_u|B_u)^2\le\abs C\,e^{-cN}.
\]
Under the hypotheses of Theorem~\ref{thm:general}, use that theorem instead, with $c_0$.

(b) Let $F$ be the computed polynomial; it is homogeneous of degree $d=N-1\ge2$ with squarefree monomials. Remove gates that compute zero and gates that no longer reach the output. Every remaining gate then computes a nonzero homogeneous multilinear polynomial with nonnegative coefficients, and
\[
  F=\sum_{j=1}^{r}g_jh_j,\qquad r\le s,
\]
where $g_j,h_j$ are nonzero homogeneous polynomials with nonnegative coefficients and $d/3<\deg g_j\le2d/3$. These are the cleaning and degree-balanced decomposition lemmas of~\cite[Lemmas~7.3 and~7.4]{Blaschke26Dimers}, whose proofs use only that every monomial of $F$ is squarefree of degree $d$; the method goes back to~\cite{Valiant80,JerrumSnir82}. There is no cancellation, so every monomial of $g_jh_j$ is the monomial of a spanning tree. If a variable $z_e$ occurred both in a monomial of $g_j$ and in a monomial of $h_j$, their product would be a monomial of $F$ containing $z_e^2$. So the variables of $g_j$ and of $h_j$ are disjoint. Let $A_j$ contain the variables of $g_j$ and none of $h_j$, and $B_j:=E\setminus A_j$. The trees whose monomials occur in $g_jh_j$ form a rectangle $R_j$ across $A_j|B_j$, and both $A_j$ and $B_j$ contain at least $d/3=(N-1)/3$ edges, since a monomial of $g_j$ has more than $d/3$ edges in $A_j$ and a monomial of $h_j$ has at least $d/3$ edges in $B_j$. Theorem~\ref{thm:touch} with $m_0=(N-1)/3\ge2N/9$ bounds each $\beta(\UST(G);A_j|B_j)^2$ by $e^{-2\kappa N/9}=e^{-2cN/3}$. Every monomial of $F$ occurs in some $g_jh_j$, so the $R_j$ cover $\cT(G)$, and Corollary~\ref{cor:rectangle} gives $1\le\sum_j\UST(G)(R_j)\le s\,e^{-2cN/3}$. Under the hypotheses of Theorem~\ref{thm:general}, Lemma~\ref{lem:mixed} with $m_0=(N-1)/3$ and Theorem~\ref{thm:cut} give the bound $e^{-c_1N}$ instead.
\end{proof}

The coefficients play no role beyond being positive. The fan-in restriction in (b) can be dropped by the grouping of Remark~\ref{rem:nodes} with degrees in place of scopes: at a high product gate whose children are all low, multiply children until the degree first exceeds $d/3$, and label the group by that gate, so again $r\le s$. For prisms and squares of expanders, $c_0$ and $c_1$ are positive constants depending only on $\Delta$ and $h_0$, and much larger than $c(D,h)$. Statement (b) is an exact monotone lower bound of the same kind as those of~\cite{CDM20Negations,CDGM22}, but for the undirected polynomial and for every 2-connected bounded-degree expander; see Section~\ref{sec:related}.

\section{Limits of the method and two variants}\label{sec:limits}

\subsection{Why some cycle condition is needed}

All our bounds go through $\max_{A|B}\beta(P;A|B)$ over balanced splits. Pendant edges make this quantity useless.

\begin{proposition}[pendant edges]\label{prop:pendant}
Let $H$ be a connected graph with $v\ge2$ vertices, at most $2v$ edges, maximum degree $\Delta$ and $h(H)\ge h_0>0$. Let $G$ be obtained by attaching a new vertex of degree one to every vertex of $H$. Then $G$ has maximum degree $\Delta+1$ and $h(G)\ge\min(1,h_0)/3$. Let $A$ be the set of pendant edges and $B=E(H)$. Then $A|B$ is balanced and $\beta(\UST(G);A|B)=1$.
\end{proposition}

\begin{proof}
Every spanning tree contains every pendant edge, so $\UST(G)$ is the product of the point mass on $x_A=\one$ and the law of $\UST(H)$ on $x_B$, and $\beta=1$. The split is balanced because $\abs A=v$, $\abs B=\abs{E(H)}$, $v-1\le\abs{E(H)}\le2v$ and $v\ge2$. For the expansion, let $S$ consist of $a$ vertices of $H$ and $b$ pendant vertices, with $0<a+b\le v$. A pendant edge is in $\partial S$ when exactly one of its endpoints is in $S$; there are at least $\abs{a-b}$ such edges. If $a\le v/2$, then $\abs{\partial S}\ge h_0a+\abs{a-b}$, which is at least $h_0(a+b)/3$ if $b\le2a$ and at least $b-a\ge(a+b)/3$ if $b>2a$. If $a>v/2$, then $v-a\ge b$, so $\abs{\partial S}\ge h_0(v-a)+(a-b)\ge h_0b+(a-b)$. This is at least $(a+b)/3$ if $a\ge2b$ and at least $h_0(a+b)/3$ if $a<2b$.
\end{proof}

For a cubic expander $H$ this gives a bounded-degree expander $G$ on which the method proves nothing. This does not show that $\UST(G)$ has small circuits. It shows that a lower bound for $G$ is at least as hard as one for $H$: a circuit for $\UST(H)$, multiplied by the indicator leaves $[x_e=1]$ of the pendant edges, is a circuit for $\UST(G)$ with $O(v)$ extra size and the same total-variation error. If $H$ has large girth, the question for $H$ is open (Section~\ref{sec:open}). The linking condition of Definition~\ref{def:linked} excludes bridges (Lemma~\ref{lem:linked-basic}), and more generally forces every edge onto a short cycle.

\subsection{Frozen regions and large girth}

Lemma~\ref{lem:local} conditions on everything outside a region, uniformly over the outside configuration. The next observation shows that the region must contain a cycle.

\begin{proposition}[acyclic regions freeze]\label{prop:frozen}
Let $I\subseteq E$ be acyclic, and let $T_0$ be a spanning tree with $I\subseteq T_0$. Then, conditionally on $T\setminus I=T_0\setminus I$, we have $T=T_0$ almost surely.
\end{proposition}

\begin{proof}
A completion $J\subseteq I$ must satisfy $\abs J=N-1-\abs{T_0\setminus I}=\abs I$, so $J=I$.
\end{proof}

So a bound like Lemma~\ref{lem:local}, which must hold for every outside configuration, needs a region that contains a cycle. In a graph of girth $g$, any such region spans at least $g$ vertices, and at most $N/g$ vertex-disjoint cycles exist. With the regions $I(C)$ used here, Theorem~\ref{thm:cut} gives an exponent of at most about $(N/g)D^{-2(g-1)}$ in a $D$-regular graph. This tends to zero once $g\ge(\frac12+\eta)\log_DN$ for a fixed $\eta>0$, which is the regime of the best explicit expanders~\cite{LPS88}. Random regular graphs have very few short cycles, so the method gives nothing there either. This is a limit of our local lemma, not a lower bound on circuit size. A different choice of regions, or bounds that hold only for typical outside configurations, might still work.

\subsection{Unbounded degree}

On the complete graph $K_n$, Lemma~\ref{lem:local} gives $p=(n-1)^{-2}$ for triangles, and at most $n/3$ vertex-disjoint triangles exist. The resulting exponent $rp^2$ is $O(n^{-3})$, and longer cycles are worse, so the method as it stands gives nothing. The known exponential lower bounds for $\UST(K_n)$~\cite{MartensMedabalimi14,ZhangHoltzenVdB20} are for exact representation; see Section~\ref{sec:related}.

\subsection{What is and is not used}

The proof uses three properties of the uniform spanning tree: a forbidden local pattern (no cycle), the spatial Markov property in the form of Lemma~\ref{lem:contract}, and the fact that, given everything outside the edges incident to a short cycle, each path obtained by deleting one edge of the cycle keeps probability bounded below. Nothing determinantal is used; the determinant appears only in the upper bound. The transfer-current theorem of Burton and Pemantle~\cite{BurtonPemantle93}, which makes the edge set of the uniform spanning tree a determinantal process, plays no role.

We do not know how to run this argument for uniform perfect matchings. There, conditioning on the outside of a small region can force the matching inside it. For example, in the grid, if the outside leaves unmatched exactly two vertices next to a unit square, and these are adjacent to two neighbouring corners of the square, the completion is unique. The dimer paper~\cite{Blaschke26Dimers} used a global argument with kernel estimates instead. Its bound $e^{\Theta(\sqrt n)}$ for $n$ edges is $e^{\Theta(L)}$ on the $L\times L$ grid, the same order as Corollary~\ref{cor:grid}.

\begin{remark}[weighted spanning trees]\label{rem:weighted}
Let $P_w(T)\propto\prod_{e\in T}w_e$ with $w_e\in[w_{\min},w_{\max}]$, $w_{\min}>0$. In the proof of Lemma~\ref{lem:local}, the numerator is at least $w_{\min}^{\abs C-1}\prod_jW_j$, where $W_j$ is the total weight of the edges at $K_j$, and the parent-edge encoding bounds the denominator by $\prod_jW_j\prod_{y\ne x_0}\sum_{e\ni y}w_e$. So the local bound becomes $(w_{\min}/(Dw_{\max}))^{\abs C-1}$, and Theorems~\ref{thm:cut} and~\ref{thm:general} hold with this $p$.
\end{remark}

\subsection{Uniform spanning forests}

Let $\USF(G)$ be the uniform distribution over \emph{all} spanning forests of $G$, that is, all acyclic subsets of $E$, with any number of components. (We do not mean rooted forests, forests with a prescribed number of components, or the uniform spanning forest measures of infinite graphs.) It also forbids cycles, but Lemma~\ref{lem:local} does not carry over unchanged. On a triangle there are seven forests, so for each edge $e$ the probability that the forest is $E(C)\setminus\{e\}$ is $1/7$, below the value $1/4=D^{-(\abs C-1)}$ of Lemma~\ref{lem:local}. A different local bound holds.

\begin{proposition}[forest local bound]\label{prop:forest}
Let $\Phi\sim\USF(G)$, let $C$ be a cycle, $e\in E(C)$ and $x_0\in V(C)$, and let $F\subseteq E\setminus I(C)$ be a forest. Then
\[
  \Pr\bigl[\,\Phi\cap E(C)=E(C)\setminus\{e\}\ \bigm|\ \Phi\setminus I(C)=F\,\bigr]\ \ge\ \prod_{y\in V(C)\setminus\{x_0\}}\frac1{\deg(y)+1}\ \ge\ (D+1)^{-(\abs C-1)} .
\]
Consequently Theorem~\ref{thm:cut} holds for $\USF(G)$ with $p_k$ replaced by $\max_{x_0}\prod_{y\ne x_0}(\deg(y)+1)^{-1}$, and Theorem~\ref{thm:general} holds for $\USF(G)$ with $p=(D+1)^{-(\ell-1)}$.
\end{proposition}

\begin{proof}
Define $K_1,\dots,K_s$ and $\mathcal M_F$ as in Lemma~\ref{lem:contract}, and let $d_j\ge0$ be the degree of $K_j$ in $\mathcal M_F$. For $J\subseteq I(C)$, the set $F\cup J$ is acyclic if and only if $J$ is acyclic in $\mathcal M_F$. Indeed, a cycle of $F\cup J$ must use edges of $J$, since $F$ is acyclic, and contracting the components of $F$ maps it to a closed trail in $\mathcal M_F$ that uses each of its $J$-edges once, which contains a cycle. Conversely, a cycle of $\mathcal M_F$ lifts to a cycle of $F\cup J$ by joining consecutive edges through the connected components of $F$. Since $\Phi$ is uniform, the conditional law of $\Phi\cap I(C)$ is uniform over the forests of $\mathcal M_F$.

\emph{Numerator.} Take the path $E(C)\setminus\{e\}$, and for each $K_j$ either no edge or one of its $d_j$ edges. Each such set is a forest, so there are at least $\prod_j(d_j+1)$ forests with $J\cap E(C)=E(C)\setminus\{e\}$.

\emph{Denominator.} Root each component of a forest $J$ of $\mathcal M_F$ at $x_0$ if it contains $x_0$, and at any fixed node otherwise. Record for each node other than $x_0$ its parent edge, or the symbol ``root''. The forest $J$ is the set of recorded edges, so the encoding is injective. The node $K_j$ has at most $d_j+1$ possible records and $y\in V(C)\setminus\{x_0\}$ has at most $\deg(y)+1$. So there are at most $\prod_j(d_j+1)\prod_{y\ne x_0}(\deg(y)+1)$ forests.

The rest of the proof of Theorem~\ref{thm:cut} uses only that $\Phi$ contains no cycle, that the past is a function of $\Phi\setminus I(C_k)$, and the local bound.
\end{proof}

This bound has been checked exactly on small graphs (Appendix~\ref{app:checks}). It gives a $2^{\Omega(N)}$ lower bound for approximating $\USF(G)$ on prisms and squares of expanders, but this is not a separation. Counting spanning forests is \#P-hard for general graphs~\cite{JVW90}, and we know of no polynomial-size generating circuit for $\USF$ on these families.

\section{Related work}\label{sec:related}

Our literature search was targeted, not exhaustive. The comparisons below refer to the statements we checked in the cited sources.

\subsection{Spanning trees and determinantal distributions in probabilistic circuits}

Martens and Medabalimi~\cite{MartensMedabalimi14} prove that the uniform distribution on spanning trees of the complete graph $K_m$ cannot be represented by decomposable and complete (smooth) sum-product networks of size below $2^{m/30240}$ for $m\ge20$, and the same holds for approximation ``arbitrarily well'' by networks of bounded size (their Theorem~37). Their proof writes the network as a sum of $s^2$ ``weak'' functions and shows that each is nonzero on only an exponentially small fraction of the trees; it is a support argument. It does not give a bound at a fixed total-variation error, and the exponent is the number of vertices, which is the square root of the number of variables. Zhang, Holtzen and Van den Broeck~\cite{ZhangHoltzenVdB20} restate this result (their Theorem~5), observe that the uniform spanning tree of $K_n$ is a determinantal point process (Theorem~4), and prove that some determinantal point processes have no polynomial-size deterministic PC. Zhang, Juba and Van den Broeck~\cite{ZhangJubaVdB21} introduce probabilistic generating circuits, which allow subtraction, show that they subsume determinantal point processes, and cite the two results above for the corresponding separation from decomposable PCs. Agarwal and Bl\"aser~\cite{AgarwalBlaser24} and Broadrick, Zhang and Van den Broeck~\cite{BroadrickZhangVdB24} study the power of such circuits further.

Theorem~\ref{thm:main} instead gives a fixed-error bound without assuming determinism or exact representation. It holds for sparse bounded-degree graphs, with an exponent linear in the number of variables. Section~\ref{sec:limits} explains why our method says nothing about $K_n$.

\subsection{Constant-error lower bounds for probabilistic circuits}

Leland and Choi~\cite{LelandChoi25} prove that a \emph{deterministic} decomposable PC that approximates a certain distribution within total variation $\eps\le\frac1{16}-\Omega(1/\mathrm{poly}(n^2))$ has size $2^{\Omega(n)}$ (their Theorem~4.1). The target is built from Sauerhoff's function~\cite{Sauerhoff03} and is itself represented by a polynomial-size nondeterministic decomposable PC. Their proof prunes inputs of small probability, which uses determinism, and reduces to weak approximation by d-DNNFs. Our lower bound does not assume determinism, and our target has a polynomial-size generating circuit with subtraction rather than a small nonnegative one.

The dimer paper~\cite{Blaschke26Dimers} proves an $\exp(\Theta(\sqrt n))$ constant-error bound for uniform domino tilings of the square grid with $n$ edges, and its Section~10.4 notes a stronger generic separation: the uniform distribution on a suitable random linear code, which has easy inference by Gaussian elimination, needs $\exp(\Omega(n))$ size at constant error, by the rank bounds of de Colnet and Mengel~\cite{DeColnetMengel20}. That example is generic rather than natural, and its tractability comes from linear algebra over $\mathbb F_2$. De Colnet and Mengel study weak and strong approximation in knowledge compilation, where a Boolean function is approximated by another Boolean function. Their results do not give ours as a black-box corollary. A distribution within small total variation of $P$ can have full support (take $(1-\delta)P+\delta U$ with $U$ uniform), so approximating $P$ says nothing about approximating its support as a set. Their robust transfer also uses covers by disjoint rectangles, which is why it targets deterministic circuits, while Lemma~\ref{lem:overlap} allows overlapping terms with different splits.

Squared and subtractive circuits~\cite{LoconteEtAl24,WangVdB25} are another way to go beyond nonnegative PCs; we do not consider them.

\subsection{Knowledge compilation and communication complexity}

Bova, Capelli, Mengel and Slivovsky~\cite[Thm.~6]{BovaCapelliMengelSlivovsky16} show that a DNNF of size $s$ yields a cover of its models by at most $s$ balanced rectangles, and use multi-partition communication complexity~\cite{DurisEtAl04} to prove DNNF lower bounds. Sauerhoff~\cite{Sauerhoff03} studies approximation of Boolean functions by such rectangles. Corollary~\ref{cor:support} is a rectangle-cover argument of this kind, with the rectangle bound coming from Corollary~\ref{cor:rectangle}.

The information-complexity method of Bar-Yossef, Jayram, Kumar and Sivakumar~\cite{BJKS04} uses Hellinger distance, the Bhattacharyya coefficient for products of distributions, and the cut-and-paste property of protocols. See also later work such as Jayram's~\cite{Jayram09}. Braun and Pokutta~\cite{BraunPokutta16} use Hellinger-type arguments for robust nonnegative-rank bounds for unique disjointness. Our chain rule (Lemma~\ref{lem:chain}) is a case of Exercise~I.42(b) of Polyanskiy and Wu~\cite{PolyanskiyWu25}, and the identification of the optimized affinity with a matrix norm is Lemma~6 of Lapidoth and Pfister~\cite{LapidothPfister19}. The forbidden outcome $u_i=v_i=1$ of our cycle bits has the form of set disjointness, which connects the argument to the Hellinger methods for disjointness and AND~\cite{BJKS04,Jayram09}. We did not find the combination with a spatial Markov property that we use here, but given how short the argument is, we would not be surprised if similar arguments appear elsewhere.

\subsection{Monotone arithmetic circuits}

Jerrum and Snir~\cite{JerrumSnir82} proved $2^{\Omega(n)}$ monotone lower bounds for the directed spanning-tree polynomial of the complete graph, which has $\Theta(n^2)$ variables. Chattopadhyay, Datta, Ghosal and Mukhopadhyay~\cite[Thm.~1.1]{CDGM22} prove that for a sufficiently large constant $d$ and every $d$-regular expander with second eigenvalue at most $d^{1-\epsilon}$, every monotone circuit for the directed spanning-tree polynomial (one variable per arc, trees directed towards a root) has size $2^{\Omega(n)}$. This is the closest exact analogue of our result that we found. It differs in three ways: it is exact rather than approximate; it concerns the directed polynomial, from which the undirected one is obtained by identifying the two arc variables of each edge (so lower bounds do not transfer from directed to undirected); and it needs large degree, while our graphs have any fixed degree but must be linked. Related monotone lower bounds via communication complexity are in~\cite{ChattopadhyayDattaMukhopadhyay21}. Since the undirected polynomial is a substitution instance of the directed one, a monotone lower bound for the undirected polynomial implies one for the directed polynomial of the same graph, but not conversely. So Corollary~\ref{cor:support}(b) also gives directed bounds for our families. For the undirected polynomial, Jukna and Seiwert~\cite[Thm.~2]{JuknaSeiwert19} prove a $2^{\Omega(\sqrt n)}$ monotone arithmetic lower bound on the complete graph $K_n$, and the exact uniform-spanning-tree lower bound of Martens and Medabalimi~\cite[Thm.~37]{MartensMedabalimi14} also transfers to monotone circuits for the complete-graph polynomial, since replacing $z_e$ by the leaf $[x_e=1]$ and padding with $[x_e=0]$ turns such a circuit into a smooth decomposable PC with polynomial overhead. So exponential exact lower bounds for undirected spanning-tree polynomials are not new in general. What Corollary~\ref{cor:support}(b) adds is bounded-degree families with an exponent linear in the number of variables. A targeted search found no earlier result of this form, but we cannot rule one out.

\subsection{Product approximations of uniform spanning trees}

M\'esz\'aros and Peled~\cite{MeszarosPeled26} prove that for a graph of minimum degree $k\ge2$ there is a coupling of the uniform spanning tree with the random $1$-out subgraph (each vertex picks one incident edge independently) whose expected number of differing edges is $O(n\log(k)/\sqrt k)$. They prove analogous couplings for a wide class of determinantal processes. This does not conflict with our result, though the $1$-out subgraph is close to a product distribution. Their bound is on the expected number of edits, a Wasserstein-type distance, not on total variation; a small expected edit distance per vertex is compatible with total variation close to one. The $1$-out law is a product over vertices, not over a split of the edges, since two vertices can choose the same edge. And at fixed degree the bound is a fixed fraction of $n$, which is small only when $k$ is large. Our result is for fixed degree and total variation.

\subsection{Expanders}

For background on expanders and explicit constructions see Hoory, Linial and Wigderson~\cite{HLW06} and Lubotzky, Phillips and Sarnak~\cite{LPS88}. Our geometric lemma (Lemma~\ref{lem:mixed}) is a giant-component argument of a standard kind.

\section{Open questions}\label{sec:open}

\begin{enumerate}
\item \emph{Large girth.} Does the uniform spanning tree of a bounded-degree expander of girth $\Omega(\log N)$, such as a Ramanujan graph~\cite{LPS88} or a random regular graph, need $2^{\Omega(N)}$-size nonnegative decomposable PCs at constant total-variation error? Our local lemma needs short cycles (Section~\ref{sec:limits}); a proof here would need a different source of conditional non-degeneracy, perhaps along long cycles with bounds that hold only for typical outside configurations.

\item \emph{Unbounded degree.} Is there a constant-error exponential lower bound for $\UST(K_n)$, with $n$ vertices and $\binom n2$ variables? Exact representation is known to need $2^{\Omega(n)}$ size~\cite{MartensMedabalimi14}. Our method gives nothing here. More generally, how does the exponent depend on the degree in families such as products or powers of expanders with growing degree?

\item \emph{Constants.} The constant in Example~\ref{ex:cubic} is about $2\cdot10^{-8}$. What is the correct exponent for, say, prisms of cubic Ramanujan graphs or for the $L\times L$ grid? Natural improvements are a sharper local lemma, using all mixed cycles rather than a vertex-disjoint subfamily, and a direct bound on $\opnorm{\sqrt\mu}$ for the joint law of all cycle bits.

\item \emph{Other models.} Does the bound extend to nonnegative circuits that are decomposable but not smooth, with a size measure independent of smoothing? To structured or deterministic variants with sharper constants? To squared or subtractive circuits~\cite{LoconteEtAl24,WangVdB25}, where the overlap lemma no longer applies?

\item \emph{Other determinantal distributions.} The proof uses only a forbidden local pattern and a spatial Markov property. Which other distributions with determinant-based inference, such as uniform bases of other regular matroids, or rooted spanning forests, whose weighted count is $\det(L+qI)$, admit the same local argument on expander-like families?
\end{enumerate}


\appendix
\section{Machine checks}\label{app:checks}

None of the proofs depends on these computations. They test the lemmas on small graphs, recompute the constants in exact arithmetic, and show how loose the bounds are. Small graphs say nothing about rates in $N$, and the computations are not evidence for the asymptotic statements.

The scripts are in the \texttt{verify/} folder of \url{https://github.com/chrisb4096-alt/expander-spanning-trees}, with their output in \texttt{verify/out/}. The command \texttt{verify/run\_all.sh} compiles the C++ helper with \texttt{g++}, runs every script with \texttt{uv run}, saves the output, and stops at the first failure. Every script exits with a nonzero status if a checked inequality fails, so a run that prints \texttt{ALL CHECKS PASSED} has checked everything listed below. The Python scripts use NumPy, SciPy and NetworkX where noted, and otherwise only the standard library.

\begin{description}
\item[\texttt{constants.py}] Corollary~\ref{cor:constants} and Example~\ref{ex:cubic} in exact rational arithmetic. For the cubic prism with $h(H)\ge1/5$ it checks $t\ge N/960$ (the other term is $N/732$), $r_0\ge N/11520$ and $r_0p^2\ge N/47185920$, and it prints the per-cycle factors $0.9998779222$ and $0.9998778924$ and the exponent $2.119535\cdot10^{-8}$. It checks that the closed forms of Corollary~\ref{cor:constants} lie below the exact value of $r_0p^2/N$ at $77$ pairs $(\Delta,h_0)$ with $2\le\Delta\le12$.

\item[\texttt{check\_tensorization.py}] Lemmas~\ref{lem:chain} and~\ref{lem:two-by-two} (NumPy). It builds $6000$ random laws on $(\{0,1\}^2)^r$ by the chain rule, with conditionals that depend on the whole past, have no mass on $(1,1)$ and mass at least $p$ on $(0,1)$ and $(1,0)$, including the extreme points of this polytope. It checks $\opnorm{\sqrt\mu}\le\prod_k\gamma_k\le((1+\sqrt{1-4p^2})/2)^{r/2}\le(1-p^2)^{r/2}$ up to floating-point tolerance $10^{-12}$.

\item[\texttt{check\_chain.py} and \texttt{enum\_trees.cpp}] The whole chain of Theorem~\ref{thm:cut} on five graphs: the prisms over $K_4$, $C_6$ and $K_{3,3}$, the square of $C_{10}$, and the $4\times4$ grid, with $3456$ to $405000$ spanning trees, enumerated by the C++ helper (NumPy, SciPy). For each graph it first checks Lemma~\ref{lem:local} for every cycle of length at most $4$, every edge on it and every outside configuration; the smallest ratio of the true conditional probability to the bound is between $1.92$ and $2.67$. It then takes five balanced splits per graph, three random and two found by local search to make $\beta$ large, and checks
\[
  \beta(\UST(G);A|B)\le\opnorm{\sqrt\mu}\le\prod_k\gamma_k\le\prod_k(1-p_k^2)^{1/2}
\]
for a greedily chosen family of disjoint mixed cycles. On these $25$ instances the geometric mean of the $\gamma_k$ is between $0.94$ and $0.99$, while the bound per cycle is between $0.998$ and $0.99988$.

\item[\texttt{check\_forest.py}] Lemma~\ref{lem:local} and Proposition~\ref{prop:forest} in exact integer arithmetic on eight small graphs with $3$ to $10$ vertices, for every cycle of length at most $4$, every edge on it and every outside configuration. It also checks the triangle: $7$ forests, conditional probability $1/7$, below the tree bound $1/4$ and above the forest bound $1/9$.

\item[\texttt{independent\_check.py}] A separate implementation written for review (NumPy, NetworkX), sharing no code with the other scripts. It checks the chain inequality on $2805$ cases ($405$ exhaustive depth-two extreme-point assignments and $2400$ random chains); Lemma~\ref{lem:local} for cycles of length at most $4$ on all connected graphs with $3$ to $5$ vertices and four further graphs ($3858$ conditional inequalities, smallest ratio to the bound $4/3$); the chain of Theorem~\ref{thm:cut} for single cycles and ordered pairs of disjoint cycles on $2894$ balanced splits; and the geometric facts behind Lemmas~\ref{lem:mixed}, \ref{lem:prism} and~\ref{lem:square}, namely the expansion, linking and $\kappa$ bounds and the count of mixed cycles, for the prisms and squares of all connected graphs with $3$ or $4$ vertices and of $C_5$ ($18$ graphs, $103878$ balanced splits). Its line about the forest counterexample records a known value and is not a computation; \texttt{check\_forest.py} computes it.
\end{description}

The local lemma is loose by a factor of about $2$ on these graphs, and the chain bound by much more: the proved factor per cycle, $1-O(D^{-2(\ell-1)})$, is far closer to $1$ than the observed one. This is the main reason for the small constants in Section~\ref{sec:main}.


\bibliographystyle{alphaurl}
\bibliography{refs}

\end{document}
